\chapter{Homologia przestrzeni}
\label{ch:homologia-przestrzeni}

W Rozdziale~\ref{ch:algebra-homologiczna} homologia była
niezmiennikiem \emph{kompleksu łańcuchowego}.
Teraz przypiszemy taki kompleks przestrzeni. Najpierw
użyjemy sympleksów z Rozdziału~\ref{ch:topologia-algebraiczna-podstawy};
potem zastąpimy je wszystkimi ciągłymi obrazami sympleksów,
aby móc badać przestrzenie bez wybranej triangulacji. Na końcu
zobaczymy, dlaczego do obliczeń zwykle wystarczy
kilka komórek CW.

Intuicja jest jedna: cykl to suma zorientowanych kawałków
bez brzegu, a brzeg to cykl, który ogranicza kawałek
o wymiar wyższy. Iloraz
\[
 H_n(X;\Z)=
 \frac{\{\text{$n$-cykle w $X$}\}}
 {\{\text{brzegi $(n+1)$-łańcuchów w $X$}\}}
\]
odróżnia na przykład pętlę będącą brzegiem na tarczy od pętli
obiegającej dziurę. Nie chodzi przy tym o sam kształt pętli:
ważne jest, jakie wypełnienia należą do badanej przestrzeni.
Pracujemy nad $\Z$; zmianę
współczynników omawia
Twierdzenie~\ref{thm:homologia-char-zero},
a konkretny rachunek znajduje się poniżej.

\section{Łańcuchy symplicjalne: orientacja i brzeg}
\label{sec:homologia-symplicjalna}

W sympleksie geometrycznym kolejność wierzchołków nie zmienia
zbioru punktów, ale pozwala określić \emph{orientację}.
Dwie kolejności różniące się parzystą permutacją nadają tę samą
orientację; nieparzysta permutacja odwraca ją. Łańcuch nie jest
podzbiorem $|K|$: to formalna suma zorientowanych sympleksów
z całkowitymi krotnościami. Współczynnik $-1$ oznacza ten sam
sympleks z przeciwną orientacją, a współczynnik $2$ oznacza jego
dwukrotne wystąpienie. Dzięki temu wspólne ściany mogą się znosić.

\begin{definicja}[Kompleks łańcuchów symplicjalnych]
\index{kompleks lancuchow symplicjalnych@Kompleks łańcuchów symplicjalnych}
\label{def:homologia-symplicjalna}
Niech $K$ będzie skończonym kompleksem symplicjalnym
z Definicji~\ref{def:kompleks-symplicjalny}.
Ustalmy porządek jego wierzchołków. Grupa
$C_n^\Delta(K)$ jest wolną grupą abelową, której bazą są
$n$-sympleksy o wierzchołkach zapisanych rosnąco:
$[v_0,\ldots,v_n]$. Jeśli wierzchołki wypiszemy
w innej kolejności, przyjmujemy zasadę
$[v_{\pi(0)},\ldots,v_{\pi(n)}]
=\operatorname{sgn}(\pi)[v_0,\ldots,v_n]$.
Zatem odwrócenie orientacji zmienia znak, a dwa
różne sympleksy nadal dają niezależne generatory.
Dla $n\geq1$ określamy
\begin{equation}\label{eq:brzeg-symplicjalny}
 \partial_n[v_0,\ldots,v_n]
 =\sum_{i=0}^n(-1)^i
 [v_0,\ldots,\widehat v_i,\ldots,v_n],
 \qquad \partial_0=0.
\end{equation}
Rozszerzamy $\partial$ liniowo na sumy
o całkowitych współczynnikach. Dla $n<0$ przyjmujemy $C_n^\Delta(K)=0$.
Symbol $\widehat v_i$ oznacza pominięcie wierzchołka $v_i$;
pozostałe zachowują swój porządek.
\end{definicja}

\begin{stwierdzenie}[Brzeg respektuje orientację]
\label{prop:brzeg-orientacja}
Wzór~\eqref{eq:brzeg-symplicjalny} nie zależy od wyboru zapisu
orientacji sympleksu: jeśli zmienimy jego orientację, brzeg zmieni
znak.
\end{stwierdzenie}
\begin{proof}
Wystarczy zamienić dwa sąsiednie wierzchołki $v_i,v_{i+1}$,
bo każdą permutację można złożyć z takich zamian. Każda ściana,
w której oba te wierzchołki pozostają, zmienia orientację, więc
jej wkład w brzeg ma przeciwny znak. Dwie pozostałe ściany
powstają przez usunięcie jednego z zamienionych wierzchołków.
Przed zamianą ich suma ma postać
$(-1)^i[\ldots,v_{i+1},\ldots]
+(-1)^{i+1}[\ldots,v_i,\ldots]$;
po zamianie te same ściany otrzymują odpowiednio współczynniki
$(-1)^{i+1}$ i $(-1)^i$. Każdy wkład zmienia więc znak.
Tak samo dzieje się z całym brzegiem. Dla $0$-sympleksu brzeg
jest zerowy i nie ma nietrywialnej permutacji wierzchołków.
\end{proof}

Znak ściany nie jest ozdobą: kierunek odziedziczony
przez brzeg wypełnionego trójkąta musi być spójny.
Dla $e=[v_0,v_1]$ i $t=[v_0,v_1,v_2]$ dostajemy
\[
 \partial e=[v_1]-[v_0],\qquad
 \partial t=[v_1,v_2]-[v_0,v_2]+[v_0,v_1].
\]
Drugą równość można zapisać
$\partial t=[v_0,v_1]+[v_1,v_2]+[v_2,v_0]$.
Jest to obieg brzegu trójkąta.

\begin{figure}[!htbp]
\centering
\begin{tikzpicture}[>=Stealth,line cap=round,line join=round]
 \coordinate (a) at (0,0); \coordinate (b) at (3.3,0);
 \coordinate (c) at (1.35,2.55);
 \fill[manifoldblue!13] (a)--(b)--(c)--cycle;
 \draw[manifoldblue!65!black,line width=1.6pt]
  (a)--(b)--(c)--cycle;
 \draw[->,manifoldred,line width=1.5pt]
  (.7,-.13)--(2.5,-.13);
 \draw[->,manifoldred,line width=1.5pt]
  (3.1,.38)--(1.8,2.07);
 \draw[->,manifoldred,line width=1.5pt]
  (1.04,2.1)--(.13,.39);
 \foreach \p in {a,b,c}{\fill[manifoldgold!85!black] (\p) circle(2.8pt);}
 \node[below left] at (a) {$v_0$};
 \node[below right] at (b) {$v_1$};
 \node[above] at (c) {$v_2$};
 \node[manifoldblue!75!black] at (1.4,.84) {$t$};
 \node[align=left,anchor=west,text width=5.5cm] at (4.1,1.38)
  {$\partial t=e_{01}+e_{12}+e_{20}$\\[3pt]
   zewnętrzna pętla jest brzegiem, gdy wnętrze $t$
   należy do przestrzeni};
\end{tikzpicture}
\caption{Zorientowany brzeg trójkąta. Usunięcie jego
wnętrza pozostawia cykl, który w samym okręgu
nie jest już brzegiem żadnego $2$-łańcucha.}
\label{fig:homologia-trojkat}
\end{figure}

\begin{twierdzenie}[Brzeg brzegu znika]
\label{thm:brzeg-brzegu}
Odwzorowania~\eqref{eq:brzeg-symplicjalny}
spełniają $\partial_{n-1}\circ\partial_n=0$.
Wobec tego
\[
 H_n^\Delta(K):=
 \ker(\partial_n:C_n^\Delta(K)\to C_{n-1}^\Delta(K))
 \big/
 \operatorname{im}(\partial_{n+1})
\]
jest dobrze określoną grupą. Licznik składa
się z \emph{cykli}, mianownik z \emph{brzegów}.
\end{twierdzenie}
\begin{proof}
Dla $n=1$ wynika to z $\partial_0=0$.
Niech $n\geq2$ i weźmy generator
$[v_0,\ldots,v_n]$. Każda jego ściana
o wymiarze $n-2$ powstaje po usunięciu
pewnych dwóch wierzchołków $v_i,v_j$,
$i<j$, dokładnie na dwa sposoby.
Usunięcie najpierw $v_i$, a następnie
$v_j$ (ma on już pozycję $j-1$) daje
współczynnik $(-1)^i(-1)^{j-1}$.
W przeciwnej kolejności współczynnik
wynosi $(-1)^j(-1)^i$.
Ich suma to
$(-1)^{i+j-1}+(-1)^{i+j}=0$.
Dotyczy to każdej ściany z osobna,
więc $\partial^2$ znika na bazie,
a przez liniowość na całej grupie.
Stąd $\operatorname{im}\partial_{n+1}
\subseteq\ker\partial_n$, co pozwala
utworzyć podany iloraz.
\end{proof}

Zapis $\partial^2$ użyty w dowodzie oznacza złożenie dwóch kolejnych
operatorów brzegu, czyli $\partial_{n-1}\circ\partial_n$ w stopniu $n$.

\begin{przyklad}[Odcinek, okrąg, tarcza]
\label{prz:homologia-odcinek-okrag}
Odcinek ma dwa wierzchołki $v_0,v_1$
i jedną krawędź $e$. Odwzorowanie
$\partial_1:\Z e\to\Z v_0\oplus\Z v_1$
przypisuje $me\mapsto m(v_1-v_0)$.
Ma zerowe jądro, więc $H_1^\Delta=0$,
a w ilorazie $H_0^\Delta$ obie klasy
wierzchołków są równe. Otrzymujemy $\Z$.

Teraz weźmy obwód trójkąta, bez wnętrza,
i ustawmy kierunki $e_{01},e_{12},e_{20}$.
Każdy $1$-łańcuch
$ae_{01}+be_{12}+ce_{20}$ jest cyklem
dokładnie wtedy, gdy $a=b=c$
(współczynnik przy każdym wierzchołku
w jego brzegu musi być zerem).
Po podstawieniu brzegów trzech krawędzi otrzymujemy
\[
 \begin{aligned}
 \partial(ae_{01}+be_{12}+ce_{20})
 &=a(v_1-v_0)+b(v_2-v_1)+c(v_0-v_2)\\
 &=(c-a)v_0+(a-b)v_1+(b-c)v_2.
 \end{aligned}
\]
Zerowanie tych trzech współczynników daje właśnie $a=b=c$.
Nie ma tu $2$-sympleksów, a więc
$H_1^\Delta\cong\Z$ z generatorem
$e_{01}+e_{12}+e_{20}$.
Jeśli wstawimy wnętrze $t$, ten sam
cykl staje się $\partial t$ i
$H_1^\Delta$ znika. Rysunek~\ref{fig:homologia-trojkat}
pokazuje dokładnie tę zmianę.
\end{przyklad}

\begin{przyklad}[Dwa trójkąty i jedna wspólna krawędź]
\label{prz:homologia-kwadrat}
Kwadrat o kolejnych wierzchołkach $v_0,v_1,v_2,v_3$
podzielmy przekątną $[v_0,v_2]$.
Weźmy uporządkowane bazy
$(e_{01},e_{12},e_{23},e_{03},e_{02})$ w stopniu $1$
oraz $(t_{012},t_{023})$ w stopniu $2$.
Wszystkie indeksy w tych symbolach rosną; w szczególności
$e_{03}$ biegnie od $v_0$ do $v_3$, przeciwnie do obiegu brzegu
kwadratu. Z definicji
$\partial t_{012}=e_{12}-e_{02}+e_{01}$ oraz
$\partial t_{023}=e_{23}-e_{03}+e_{02}$.
Macierze brzegu w tych bazach i bazie $(v_0,v_1,v_2,v_3)$ wynoszą
\[
 [\partial_1]=
 \begin{pmatrix}
 -1&0&0&-1&-1\\
 1&-1&0&0&0\\
 0&1&-1&0&1\\
 0&0&1&1&0
 \end{pmatrix},\qquad
 [\partial_2]=
 \begin{pmatrix}1&0\\1&0\\0&1\\0&-1\\-1&1\end{pmatrix}.
\]
Każda kolumna jest brzegiem jednego generatora, a nie listą
niezależnych relacji. Dla wektora współczynników $(a,b,c,d,e)$
równanie $\partial_1(a,b,c,d,e)=0$ daje kolejno
$b=a$, $d=-c$ i $e=c-a$; równanie przy $v_0$ jest wtedy
automatycznie spełnione. Zatem
\[
 (a,b,c,d,e)
 =a(1,1,0,0,-1)+c(0,0,1,-1,1)
 =a\,\partial t_{012}+c\,\partial t_{023}.
\]
Każdy $1$-cykl jest brzegiem, więc $H_1^\Delta=0$.
Ponadto $\partial_2$ jest iniektywny: współczynniki przy
$e_{01}$ i $e_{23}$ odzyskują oba współczynniki $2$-łańcucha.
Stąd $H_2^\Delta=0$. Relacje pochodzące z krawędzi utożsamiają
wszystkie cztery wierzchołki, pozostawiając $H_0^\Delta\cong\Z$.

Suma $t_{012}+t_{023}$ ma brzeg
$e_{01}+e_{12}+e_{23}-e_{03}$: przekątna znosi się z przeciwnym
znakiem. Gdy usuniemy tylko wnętrza obu trójkątów, pozostawiając
pięć krawędzi, grupa cykli pozostaje ta sama, lecz nie ma już
$2$-łańcuchów. Wtedy $H_1^\Delta\cong\Z^2$.
To na tych samych macierzach pokazuje różnicę między cyklem
a klasą cyklu modulo dostępne wypełnienia.
\end{przyklad}

W $\Delta$-kompleksie z
Definicji~\ref{def:delta-kompleks} postępujemy
podobnie: każdy \emph{sympleks charakterystyczny}
jest osobnym generatorem, nawet jeżeli jego
różne ściany mają ten sam obraz. Brzeg
jest sumą obrazów ścian z ich znakami;
po identyfikacji mogą się one wzajemnie
skasować. Tak właśnie sklejony odcinek
z jednym wierzchołkiem ma $\partial e=v-v=0$.

\begin{stwierdzenie}[Składowe i $H_0$]
\label{prop:h0-symplicjalne}
Jeżeli $K$ ma $c$ składowych łukowych,
to $H_0^\Delta(K)\cong\Z^c$.
\end{stwierdzenie}
\begin{proof}
Ponieważ $\partial_0=0$, grupa $H_0$
to wolna grupa na wierzchołkach
modulo relacje
$v_j-v_i=\partial[v_i,v_j]$
dla krawędzi. Idąc po drodze krawędziowej,
utożsamiamy klasy wszystkich wierzchołków
jej składowej. Dwa wierzchołki tej samej
składowej $|K|$ można połączyć taką drogą:
wszystkie wierzchołki jednego sympleksu należą do tej samej
składowej grafu krawędzi. Sympleksy przypisane jednej takiej
składowej tworzą podkompleks. Jego realizacja i jej dopełnienie
są domknięte jako skończone sumy domkniętych sympleksów,
więc są też otwarte. Droga nie może przejść między tymi zbiorami.
Każdy punkt sympleksu można natomiast połączyć odcinkiem z wierzchołkiem.
Składowe grafu odpowiadają zatem dokładnie składowym łukowym $|K|$. Wreszcie żadna
relacja nie łączy różnych składowych,
bo obie końcówki każdej krawędzi leżą
w jednej. W każdej składowej pozostaje
więc jeden wolny generator.
\end{proof}

\section{Łańcuchy singularne i niezmienniczość homotopijna}
\label{sec:homologia-singularna}

Triangulacja bywa trudna do znalezienia. W konstrukcji
singularnej sam wybór \emph{ciągłego} odwzorowania
standardowego sympleksu do przestrzeni tworzy
generator; nie żądamy różnowartościowości ani
ładnego kształtu jego obrazu.

\begin{definicja}[Homologia singularna]
\index{homologia singularna@Homologia singularna}
\label{def:homologia-singularna}
\emph{Singularny $n$-sympleks} w $X$ to dowolne
ciągłe odwzorowanie $\sigma:\Delta^n\to X$.
Niech $C_n(X)$ będzie wolną grupą abelową
na wszystkich takich odwzorowaniach; dla $n<0$ przyjmujemy $C_n(X)=0$.
W szczególności
dwa różne odwzorowania o tym samym obrazie są
różnymi generatorami. Niech
$\iota_i:\Delta^{n-1}\to\Delta^n$
będzie afiniczną inkluzją ściany
naprzeciw $i$-tego wierzchołka. Kładziemy
\begin{equation}\label{eq:singular-boundary}
 \partial\sigma=\sum_{i=0}^n(-1)^i
 \sigma\circ\iota_i\quad(n\geq1),
 \qquad \partial:C_0(X)\to0,
 \qquad H_n(X):=H_n(C_\bullet(X),\partial).
\end{equation}
\end{definicja}

Obliczenie $\partial^2=0$ jest słowo w słowo
rachunkiem dwóch usuwanych wierzchołków
z dowodu Twierdzenia~\ref{thm:brzeg-brzegu}.
Teraz pojawia się nawet singularny sympleks
stały: nie odrzucamy go na poziomie
łańcuchów, gdyż później jego wkład może
być brzegiem.

Warto oddzielić dwie konwencje orientacji. Dla singularnej drogi
$\gamma$ odwzorowanie $\bar\gamma(t)=\gamma(1-t)$ jest innym
generatorem niż $\gamma$, jeżeli te funkcje są różne;
nie narzucamy relacji $\bar\gamma=-\gamma$ w $C_1(X)$.
Mamy natomiast $\partial\bar\gamma=-\partial\gamma$.
Zależności między klasami pętli o przeciwnych kierunkach
wyprowadzimy z brzegów $2$-łańcuchów w
Sekcji~\ref{sec:obliczenia-homologii}.

\begin{figure}[!htbp]
\centering
\begin{tikzpicture}[>=Stealth,line cap=round,line join=round]
 \coordinate (a) at (0,0); \coordinate (b) at (2.45,0);
 \coordinate (c) at (.9,1.9);
 \fill[manifoldblue!12] (a)--(b)--(c)--cycle;
 \draw[manifoldblue!70!black,thick] (a)--(b)--(c)--cycle;
 \node[below left] at (a) {$v_0$};
 \node[below right] at (b) {$v_1$};
 \node[above] at (c) {$v_2$};
 \node at (1.05,.55) {$\Delta^2$};
 \draw[->,manifoldred,very thick] (2.95,.95)--(4.28,.95)
  node[midway,above] {$\sigma$};
 \fill[manifoldgreen!12]
  (4.95,.12) .. controls (5.7,-.65) and (6.6,.24) ..
  (7.75,-.08) .. controls (8.8,.3) and (8.3,1.4) ..
  (7.5,1.55) .. controls (6.5,2.2) and (5.2,1.66) ..
  (4.95,.12);
 \draw[manifoldgreen!70!black,thick]
  (4.95,.12) .. controls (5.7,-.65) and (6.6,.24) ..
  (7.75,-.08) .. controls (8.8,.3) and (8.3,1.4) ..
  (7.5,1.55) .. controls (6.5,2.2) and (5.2,1.66) ..
  (4.95,.12);
 \fill[manifoldblue!25,opacity=.7]
  (5.6,.45) .. controls (6.3,.1) and (7.3,.2) .. (7.9,.65)
  .. controls (8.05,1.1) and (7.3,1.6) .. (6.8,1.5)
  .. controls (6.1,1.7) and (5.8,1.05) .. (5.6,.45);
 \draw[manifoldred,thick,->]
  (5.6,.45) .. controls (6.3,.1) and (7.3,.2) .. (7.9,.65);
 \draw[manifoldred,thick,->]
  (7.9,.65) .. controls (8.05,1.1) and (7.3,1.6) .. (6.8,1.5);
 \draw[manifoldred,thick,->]
  (6.8,1.5) .. controls (6.1,1.7) and (5.8,1.05) .. (5.6,.45);
 \foreach \x/\y in {5.6/.45,7.9/.65,6.8/1.5}
  {\fill[manifoldred] (\x,\y) circle(1.8pt);}
 \node[font=\scriptsize,below left] at (5.6,.45) {$\sigma(v_0)$};
 \node[font=\scriptsize,right] at (7.9,.65) {$\sigma(v_1)$};
 \node[font=\scriptsize,above] at (6.8,1.5) {$\sigma(v_2)$};
 \node[manifoldgreen!70!black] at (8.05,1.87) {$X$};
 \node[manifoldblue!70!black] at (6.7,.8) {$\sigma(\Delta^2)$};
\end{tikzpicture}
\caption{Singularny sympleks jest odwzorowaniem do $X$, niekoniecznie
płaskim trójkątem. Jego brzeg liczymy
z ograniczeń odwzorowania $\sigma$ do trzech ścian.}
\label{fig:singular-simplex}
\end{figure}

\begin{stwierdzenie}[Ciągłe odwzorowania działają na homologii]
\label{prop:funktor-singularny}
Jeśli $f:X\to Y$ jest ciągłe, to
$f_\#:C_n(X)\to C_n(Y)$,
$f_\#(\sigma)=f\circ\sigma$,
jest odwzorowaniem łańcuchowym. Wyznacza homomorfizm
$f_*:H_n(X)\to H_n(Y)$, przy czym
$(g\circ f)_*=g_*\circ f_*$ oraz
$(\id_X)_*=\id_{H_n(X)}$.
\end{stwierdzenie}
\begin{proof}
Na pojedynczym generatorze
\[
 \partial(f\circ\sigma)
 =\sum_i(-1)^i f\circ\sigma\circ\iota_i
 =f_\#(\partial\sigma).
\]
Równość na sumach wynika z liniowości.
Zatem obraz cyklu jest cyklem, a obraz
brzegu brzegiem, więc $f_*$ jest
dobrze określony na ilorazie.
Równości dotyczące składania sprawdzamy
już na $\sigma$:
$(g_\#\circ f_\#)(\sigma)=g\circ f\circ\sigma$.
\end{proof}

\begin{przyklad}[Punkt]
\label{prz:homologia-punktu}
Dla $X=\{*\}$ w każdym stopniu jest
dokładnie jeden singularny sympleks
stały $c_n$, więc $C_n(X)\cong\Z$.
Jego brzeg to
$\partial c_n=(\sum_{i=0}^n(-1)^i)c_{n-1}$:
dla $n$ nieparzystego współczynnik
jest $0$, dla dodatniego parzystego $1$.
Dla $n=2r+1$ mamy $r+1$ par $1-1$, więc suma wynosi $0$;
dla $n=2r$ mamy $r$ takich par i ostatni składnik $+1$.
Kompleks ma więc postać
\[
 \cdots\xrightarrow{\;0\;}\Z
 \xrightarrow{\;1\;}\Z
 \xrightarrow{\;0\;}\Z
 \xrightarrow{\;1\;}\Z
 \xrightarrow{\;0\;}\Z\longrightarrow0,
\]
gdzie prawy skraj to $C_0$; precyzyjnie
$\partial_1=0$, $\partial_2=1$.
Stąd $H_0(*)=\Z$ i $H_n(*)=0$ dla $n>0$.
\end{przyklad}

\begin{stwierdzenie}[Składowe łukowe]
\label{prop:h0-singularne}
Grupa $H_0(X)$ jest wolna na składowych
łukowych $X$. W szczególności
dla przestrzeni niepustej łukowo spójnej
$H_0(X)\cong\Z$.
\end{stwierdzenie}
\begin{proof}
Każdy punkt $x$ jest singularnym
$0$-sympleksem, więc $C_0(X)$
jest wolną grupą na punktach.
Dla drogi $\gamma:[0,1]=\Delta^1\to X$
mamy $\partial\gamma=[\gamma(1)]-
[\gamma(0)]$. Punkty w jednej składowej
łukowej mają zatem tę samą klasę.
Z drugiej strony każdy składnik
brzegu dowolnego $1$-łańcucha
jest różnicą punktów połączonych drogą:
nie może utożsamić dwóch różnych
składowych. Sumowanie współczynników
punktów osobno w każdej składowej
daje odwrotny homomorfizm do
wolnej grupy na składowych.
\end{proof}

\begin{twierdzenie}[Niezmienniczość homotopijna]
\label{thm:homotopia-singularna}
Jeśli $f,g:X\to Y$ są homotopijne, to
$f_*=g_*:H_n(X)\to H_n(Y)$ w każdym stopniu.
Jeśli $X$ i $Y$ są homotopijnie równoważne,
ich grupy homologii są izomorficzne.
\end{twierdzenie}
\begin{proof}
Niech $F:X\times[0,1]\to Y$ będzie
homotopią, $F(x,0)=f(x)$,
$F(x,1)=g(x)$.
Graniastosłup $\Delta^n\times[0,1]$
dzielimy na $n+1$ zorientowanych
$(n+1)$-sympleksów. Dla $i=0,\ldots,n$
ich uporządkowane wierzchołki to
\[
 (e_0,0),\ldots,(e_i,0),
 (e_i,1),\ldots,(e_n,1).
\]
Niech $T_i:\Delta^{n+1}\to
\Delta^n\times[0,1]$ będzie odwzorowaniem
afinicznym posyłającym kolejne wierzchołki
właśnie na tę listę. Określmy
\[
 P_n(\sigma)=\sum_{i=0}^n(-1)^i\,
 F\circ(\sigma\times\id_{[0,1]})\circ T_i
 \in C_{n+1}(Y).
\]
Rozpiszmy brzeg sumy. Wewnętrzne
$n$-ściany, na których spotykają się
sąsiednie części graniastosłupa,
pojawiają się z przeciwnymi znakami.
Ściany boczne nad każdą ścianą
$\Delta^n$ tworzą dokładnie
$-P_{n-1}(\partial\sigma)$:
jeśli $j<i$, usuwamy dolny wierzchołek $(e_j,0)$
z $T_i$, z łącznym znakiem $(-1)^{i+j}$.
Na pryzmacie ściany naprzeciw $e_j$ punkt podziału ma indeks $i-1$,
więc znak w $-P_{n-1}\partial$ wynosi
$-(-1)^j(-1)^{i-1}=(-1)^{i+j}$.
Jeśli $j>i$, usuwamy górny wierzchołek $(e_j,1)$,
na pozycji $j+1$; oba rachunki dają $(-1)^{i+j+1}$.
Wewnętrzna ściana między $T_{i-1}$ i $T_i$ powstaje
przez usunięcie odpowiednio $(e_{i-1},1)$ i $(e_i,0)$,
ze znakami $(-1)^{i-1}(-1)^i=-1$ i
$(-1)^i(-1)^i=+1$: oba usuwane wierzchołki
mają pozycję $i$ w swoich uporządkowanych listach.
Pozostałe dwie ściany to górna
$g\circ\sigma$ ze znakiem $+1$
oraz dolna $f\circ\sigma$ ze
znakiem $-1$. Zatem
\begin{equation}\label{eq:prism-identity}
 \partial\circ P_n+P_{n-1}\circ\partial
 =g_\#-f_\#.
\end{equation}
Przyjmujemy $P_{-1}=0$. Dla $n=0$ jest to znana równość
,,koniec drogi minus początek'';
powyższe parowanie ścian dowodzi
jej w każdym wymiarze. W świetle
Twierdzenia~\ref{thm:homotopia-homologia}
oba odwzorowania indukują to samo na homologii.
Jeśli $u:X\to Y$, $v:Y\to X$
są odwrotne z dokładnością do homotopii,
to $(v\circ u)_*=\id$ i $(u\circ v)_*=\id$,
więc $u_*$ i $v_*$ są odwrotne.
\end{proof}

W szczególności przestrzeń ściągalna
ma homologię punktu. Jeśli
$A\subseteq X$ jest retraktem
deformacyjnym, to $H_n(A)\cong H_n(X)$.

\begin{przyklad}[Tożsamość pryzmatyczna dla drogi]
\label{prz:pryzmat-drogi}
Dla $\sigma:[0,1]\to X$ oznaczmy dolne wierzchołki kwadratu
przez $a_0,a_1$, a górne przez $b_0,b_1$.
Przez $[a_0,b_0,b_1]$ rozumiemy teraz singularny sympleks
otrzymany po afinicznym odwzorowaniu na ten trójkąt i zastosowaniu
$F\circ(\sigma\times\id_I)$; podobnie oznaczamy pozostałe części.
Definicja daje $P_1(\sigma)=[a_0,b_0,b_1]-[a_0,a_1,b_1]$.
Rozpisanie sześciu ścian pokazuje wszystkie znaki:
\[
 \begin{aligned}
 \partial\bigl(P_1(\sigma)\bigr)
 &=[b_0,b_1]-[a_0,b_1]+[a_0,b_0]\\
 &\quad-[a_1,b_1]+[a_0,b_1]-[a_0,a_1]\\
 &=g_\#(\sigma)-f_\#(\sigma)
   -P_0\bigl([\sigma(1)]-[\sigma(0)]\bigr).
 \end{aligned}
\]
Przekątna znosi się dosłownie jako to samo odwzorowanie
singularne, występujące raz ze znakiem dodatnim i raz z ujemnym.
Jeśli $\sigma$ jest cyklem, ostatni składnik znika: obrazy
cyklu przez $f$ i $g$ różnią się brzegiem.
\end{przyklad}

\begin{figure}[!htbp]
\centering
\begin{tikzpicture}[>=Stealth,line cap=round,line join=round]
 \coordinate (a0) at (0,0); \coordinate (a1) at (4,0);
 \coordinate (b0) at (0,2.6); \coordinate (b1) at (4,2.6);
 \fill[manifoldblue!12] (a0)--(b0)--(b1)--cycle;
 \fill[manifoldgold!15] (a0)--(a1)--(b1)--cycle;
 \draw[manifoldblue!70!black,thick] (a0)--(a1)--(b1)--(b0)--cycle;
 \draw[manifoldred,dashed,thick] (a0)--(b1);
 \draw[->,thick] (.4,0)--(3.6,0) node[midway,below] {$f\circ\sigma$};
 \draw[->,thick] (.4,2.6)--(3.6,2.6) node[midway,above] {$g\circ\sigma$};
 \draw[->,manifoldgreen!70!black,thick] (0,.35)--(0,2.25)
  node[midway,left] {$P_0(\sigma(0))$};
 \draw[->,manifoldgreen!70!black,thick] (4,.35)--(4,2.25)
  node[midway,right] {$P_0(\sigma(1))$};
 \node at (1,1.8) {$T_0\;(+)$}; \node at (3,.8) {$T_1\;(-)$};
 \node[below left] at (a0) {$a_0$}; \node[below right] at (a1) {$a_1$};
 \node[above left] at (b0) {$b_0$}; \node[above right] at (b1) {$b_1$};
\end{tikzpicture}
\caption{Dziedzina pryzmatu nad $1$-sympleksem. Podpisy krawędzi
oznaczają ich obrazy w $Y$. Suma zorientowanych trójkątów ma brzeg
równy górnej drodze minus dolna, z dwiema bocznymi poprawkami.}
\label{fig:pryzmat-homologia}
\end{figure}

\section{Pary przestrzeni i homologia zredukowana}
\label{sec:homologia-wzgledna}

\begin{definicja}[Grupy względne i zredukowane]
\index{grupy wzgledne i zredukowane@Grupy względne i zredukowane}
\label{def:homologia-wzgledna}
Dla $A\subseteq X$ inkluzja daje
podkompleks $C_\bullet(A)\subseteq C_\bullet(X)$.
Kładziemy
\[
 C_n(X,A)=C_n(X)/C_n(A),\qquad
 H_n(X,A)=H_n(C_\bullet(X,A)).
\]
\emph{Augmentacja} $\varepsilon:C_0(X)\to\Z$ sumuje
współczynniki punktów. Dołączamy stopień $C_{-1}=\Z$,
z różniczką $\partial_0=\varepsilon$; niższe grupy są zerowe.
Ponieważ $\varepsilon\circ\partial_1=0$, otrzymujemy kompleks,
którego homologie oznaczamy $\widetilde H_n(X)$.
Dla niepustego $X$ mamy
$\widetilde H_0(X)=\ker(\varepsilon:H_0(X)\to\Z)$,
$\widetilde H_n(X)=H_n(X)$ dla $n>0$ i
$\widetilde H_n(X)=0$ dla $n<0$.
Dla $X=\varnothing$ wszystkie grupy łańcuchów w stopniach
nieujemnych są zerowe, więc $\widetilde H_{-1}(\varnothing)=\Z$
i pozostałe grupy zredukowane są zerowe.
\end{definicja}

Konkretnie, dwa łańcuchy reprezentują ten sam element $C_n(X,A)$,
gdy różnią się łańcuchem leżącym w $A$. Klasa $c+C_n(A)$ jest
cyklem wtedy i tylko wtedy, gdy $\partial c\in C_{n-1}(A)$.
Dwa takie cykle $c,c'$ dają tę samą klasę homologii względnej,
gdy $c-c'=\partial b+a$ dla pewnych $b\in C_{n+1}(X)$
i $a\in C_n(A)$. Mamy więc dwie różne operacje ilorazowe:
najpierw pomijamy łańcuchy w $A$, a następnie brzegi w tym ilorazie.

Względny cykl może mieć brzeg w $A$:
to, co w $X$ jest otwartą powierzchnią,
po ,,pominięciu $A$'' staje się cyklem.
Na przykład zorientowany dysk reprezentuje klasę względną
względem swojego okręgu brzegowego.
Jeżeli $X$ ma skończoną liczbę $c>0$ składowych łukowych,
to $\widetilde H_0(X)\cong\Z^{c-1}$.
Ogólnie, po wyborze jednej składowej $P_0$, bazę jądra
augmentacji tworzą różnice $[P]-[P_0]$ dla $P\ne P_0$:
jest to suma prosta kopii $\Z$, także dla nieskończenie wielu składowych.

\begin{twierdzenie}[Długi ciąg pary]
\label{thm:les-pary}
Każda para $(X,A)$ daje naturalny
długi ciąg dokładny
\begin{equation}\label{eq:les-pary}
 \cdots\to H_n(A)\xrightarrow{i_*}H_n(X)
 \xrightarrow{j_*}H_n(X,A)
 \xrightarrow{\delta}H_{n-1}(A)
 \to H_{n-1}(X)\to\cdots.
\end{equation}
Jeśli $\partial c\in C_{n-1}(A)$,
to $\delta[c]=[\partial c]$.
\end{twierdzenie}
\begin{proof}
W każdej grupie łańcuchów mamy
krótki ciąg dokładny
$0\to C_n(A)\to C_n(X)\to C_n(X,A)\to0$.
Jest zgodny z brzegami, więc
Stwierdzenie~\ref{prop:ses-iloraz}
daje długi ciąg i wzór na $\delta$.
Sprawdźmy bezpośrednio znaczenie tego wzoru. Zmiana reprezentanta
na $c+a$, $a\in C_n(A)$, zmienia $\partial c$ o brzeg $\partial a$
w $A$. Zmiana klasy względnej o $\partial b$ nie zmienia jej brzegu,
bo $\partial^2b=0$. Zatem $\delta$ nie zależy od tych wyborów.

Dokładność w $H_n(X,A)$ też widać na łańcuchach. Jeśli
$\delta[c]=0$, to $\partial c=\partial a$ dla pewnego
$a\in C_n(A)$. Wtedy $c-a$ jest cyklem w $X$, a jego obraz
w grupie względnej jest równy $[c]$.
Odwrotnie, cykl z $X$ ma zerowy brzeg, więc jego obraz przez
$\delta$ jest zerowy. W $H_n(X)$ klasa znika względnie
dokładnie wtedy, gdy $c=\partial b+a$ z $a$ w $A$;
z $\partial c=0$ wynika wówczas $\partial a=0$, więc
$[c]=i_*[a]$. Wreszcie cykl $a$ w $A$ znika po włożeniu do $X$
dokładnie wtedy, gdy $a=\partial c$ w $X$, czyli
$[a]=\delta[c]$ dla względnego cyklu $c$ stopnia $n+1$.
Naturalność pochodzi z przemienności
odwzorowań łańcuchów dla odwzorowań
par $(X,A)\to(Y,B)$.
\end{proof}

\begin{przyklad}[Odcinek względem obu końców]
\label{prz:para-odcinek-konce}
Niech $I=[0,1]$ i $A=\{0,1\}$. Dla $k>0$ mamy $H_k(I)=0$
oraz $H_k(A)=0$, natomiast $H_0(A)=\Z[0]\oplus\Z[1]$
i $H_0(I)=\Z$. Fragment ciągu pary przyjmuje postać
\[
 0\longrightarrow H_1(I,A)\xrightarrow{\ \delta\ }
 \Z[0]\oplus\Z[1]\xrightarrow{\ (a,b)\mapsto a+b\ }
 \Z\longrightarrow H_0(I,A)\longrightarrow0.
\]
Oba końce należą do tej samej składowej $I$, stąd wzór
na środkową strzałkę. Jest ona surjektywna i ma jądro
$\Z([1]-[0])$. Dokładność daje $H_0(I,A)=0$ i
$H_1(I,A)\cong\Z$. Droga $c(t)=t$ nie jest cyklem zwykłym,
lecz jest cyklem względnym, bo $\partial c=[1]-[0]$ leży w $A$.
Ponieważ $\delta[c]=[1]-[0]$ jest generatorem jądra,
$[c]$ generuje grupę względną. Wyższe grupy względne są zerowe
z pozostałych fragmentów tego samego ciągu.

Dla porównania, jeśli zachowamy w $A$ tylko punkt $0$,
to $H_0(\{0\})\to H_0(I)$ jest izomorfizmem i wszystkie
grupy $H_k(I,\{0\})$ znikają. Dwa wyróżnione końce mierzą
przejście między nimi, jeden koniec nie daje takiej klasy.
\end{przyklad}

Jeśli $A\ne\varnothing$, możemy w tym ciągu zastąpić homologie
$A$ i $X$ zredukowanymi, pozostawiając zwykłą homologię względną.
Istotnie, w stopniu $-1$ inkluzja kompleksów augmentowanych
jest identycznością $\Z\to\Z$, więc ich iloraz ma tam zero
i jest właśnie kompleksem $C_\bullet(X,A)$.

Homotopia par (w każdym czasie obraz $A$
leży w $B$) daje te same odwzorowania na $H_n(X,A)$:
pryzmat z dowodu Twierdzenia~\ref{thm:homotopia-singularna}
przenosi $C_\bullet(A)$ do $C_{\bullet+1}(B)$,
więc wzór~\eqref{eq:prism-identity}
schodzi na kompleksy ilorazowe.

\section{Podział barycentryczny, wycinanie i sklejanie}
\label{sec:wycinanie-homologia}

Homologia singularna ma nieskończenie wiele generatorów nawet
dla prostych przestrzeni. Podział dużego sympleksu na małe
części pozwala jednak sprowadzić rachunek do lokalnych
fragmentów. Z tego mechanizmu wynika wycinanie.

\begin{lemat}[Małe łańcuchy]
\label{lem:male-lancuchy}
Niech $\mathcal U=\{U_\alpha\}$ będzie otwartym
pokryciem $X$. Podkompleks
$C_n^{\mathcal U}(X)\subseteq C_n(X)$
generowany przez sympleksy, których obrazy
leżą w jednym $U_\alpha$, ma te same
grupy homologii co $C_\bullet(X)$.
Inkluzja indukuje izomorfizmy w każdym stopniu.
\end{lemat}
\begin{proof}
Podamy konstrukcję potrzebną w obu kierunkach.
Podzielmy liniowo $\Delta^n$ na sympleksy
barycentryczne: wierzchołkami małego sympleksu
są barycentra łańcucha niepustych ścian
$F_0\subsetneq F_1\subsetneq\cdots\subsetneq F_n$,
zapisane w tej właśnie kolejności. Każdy taki łańcuch odpowiada
permutacji $\pi$ wierzchołków: $F_j=[e_{\pi(0)},\ldots,e_{\pi(j)}]$.
Niech $a_\pi:\Delta^n\to\Delta^n$ będzie afinicznym odwzorowaniem
na uporządkowany mały sympleks. Określamy
\[
 S_n(\sigma)=\sum_{\pi\in S_{n+1}}\operatorname{sgn}(\pi)
                 \,\sigma\circ a_\pi.
\]
Znak pochodzi z orientacji: dla permutacji identycznościowej
macierz kolejnych barycentrów ma dodatni wyznacznik w otoczce afinicznej,
a permutacja wierzchołków zmienia go o $\operatorname{sgn}(\pi)$.
Rzeczywiście, dla tej pierwszej permutacji
$b_j=(e_0+\cdots+e_j)/(j+1)$ oraz $b_0=e_0$, więc
$b_j-b_0=\sum_{\ell=1}^j(e_\ell-e_0)/(j+1)$.
Macierz tych kolumn w bazie $(e_1-e_0,\ldots,e_n-e_0)$
jest trójkątna, z przekątną $1/2,1/3,\ldots,1/(n+1)$.
Jej wyznacznik wynosi $1/(n+1)!>0$.
Kolejność barycentrów jest ważna: w kompleksie singularnym
odwrócenie parametryzacji nie jest z definicji ujemnym generatorem.
Ściany wewnętrzne dwóch części
występują z przeciwnymi znakami,
a ich ściany na zewnątrz dają
podział odpowiedniej ściany $\sigma$.
Dokładniej, wewnętrzną ścianę uzyskujemy, pomijając barycentrum $F_j$
dla $j<n$. Dwie odpowiadające jej permutacje zamieniają miejscami
pozycje $j,j+1$, więc mają przeciwne znaki, a pozostałe
uporządkowane barycentra są identyczne. Dla $j=n$ otrzymujemy
sympleks podziału ściany naprzeciw $e_{\pi(n)}$.
Jeśli $\rho$ jest permutacją pozostałych wierzchołków w ich
odziedziczonej kolejności, to
$(-1)^n\operatorname{sgn}(\pi)=(-1)^{\pi(n)}\operatorname{sgn}(\rho)$.
Istotnie, przeniesienie $e_{\pi(n)}$ z jego pierwotnego miejsca
na koniec wymaga $n-\pi(n)$ transpozycji, zatem
$\operatorname{sgn}(\pi)=(-1)^{n-\pi(n)}\operatorname{sgn}(\rho)$;
pomnożenie przez $(-1)^n$ daje wskazaną równość.
To właśnie znak tej ściany w $S\circ\partial$;
stąd $\partial\circ S=S\circ\partial$.

Wyjaśnijmy też, dlaczego $S$ nie zmienia
klasy homologii. Na standardowym sympleksie
oznaczmy jego identycznościowe odwzorowanie
przez $u_n$ i indukcyjnie zbudujmy
łańcuchy $t_n$ tak, aby dla $n\geq1$
\[
 \partial t_n+\sum_{i=0}^n(-1)^i(\iota_i)_\#t_{n-1}
       =S(u_n)-u_n. \tag{*}
\]
W stopniu $0$ obie strony są zerowe,
więc $t_0=0$. Jeśli mamy już konstrukcję
na wszystkich ścianach, to oznaczmy ich sumę
$b_n=\sum_i(-1)^i(\iota_i)_\#t_{n-1}$.
Łańcuch $z=S(u_n)-u_n-b_n$
jest cyklem: stosując $\partial$,
używamy $\partial\circ S=S\circ\partial$,
równości (*) na ścianach
i $\partial^2=0$.
Rozpiszmy ten krok. Po wstawieniu (*) dla $t_{n-1}$
otrzymujemy
\[
 \partial b_n
 =\sum_{i=0}^n(-1)^i(\iota_i)_\#
       \bigl(S(u_{n-1})-u_{n-1}\bigr)
 =S(\partial u_n)-\partial u_n.
\]
Dodatkowa suma zawierająca $t_{n-2}$ znika:
każda ściana kowymiaru dwa występuje dwa razy
z przeciwnymi znakami, dokładnie jak w dowodzie
Twierdzenia~\ref{thm:brzeg-brzegu}. Dla $n=1$ nie ma tej sumy,
bo $t_0=0$. Stąd
$\partial z=S(\partial u_n)-\partial u_n-\partial b_n=0$.
W wypukłym $\Delta^n$ taki cykl jest
brzegiem: z barycentrum $b$ łączymy
każdy jego singularny sympleks prostym
stożkiem. Dla $\tau:\Delta^k\to\Delta^n$ stożek $b*\tau$
ma wzór
\[
 (b*\tau)(s_0,\ldots,s_{k+1})
 =s_0b+(1-s_0)\tau\left(\frac{s_1}{1-s_0},\ldots,
                         \frac{s_{k+1}}{1-s_0}\right)
 \quad(s_0<1),
\]
z wartością $b$ dla $s_0=1$. Jest ciągły także w tym wierzchołku,
bo obraz $\tau$ jest ograniczony.
Wyprowadźmy jego brzeg dla $k\geq1$. Ściana naprzeciw
pierwszego wierzchołka $b$ ma $s_0=0$ i jest dokładnie $\tau$.
Usunięcie $i$-tego wierzchołka z części należącej do $\tau$
ma natomiast w całej liście pozycję $i+1$
i pozostawia stożek nad $i$-tą ścianą $\tau$. Zatem
\[
 \partial(b*\tau)
 =\tau+\sum_{i=0}^k(-1)^{i+1}b*(\tau\circ\iota_i)
 =\tau-b*(\partial\tau).
\]
Dla $k=0$ brzeg odcinka od $b$ do punktu $\tau$
wynosi $\tau-b$. Przedłużamy stożkowanie liniowo.
Ponieważ nasz cykl $z$ ma stopień $n\geq1$,
$\partial(b*z)=z-b*(\partial z)=z$.
Możemy więc wybrać $t_n=b*z$, co daje (*).
Przenosząc $t_n$ przez dowolne
$\sigma:\Delta^n\to X$, otrzymujemy
operator $T_n(\sigma)=\sigma_\#t_n$ o nośniku zawartym
w $\sigma(\Delta^n)$, który spełnia
\[
 \partial\circ T+T\circ\partial=S-\id.
\]
Zatem $S^r-\id=\partial\circ T^{(r)}+
T^{(r)}\circ\partial$, gdzie
$T^{(r)}=T+S\circ T+\cdots+S^{r-1}\circ T$.
Potęga $S^r$ oznacza $r$ kolejnych zastosowań operatora $S$.
Wyprowadzenie jest sumą teleskopową: ponieważ $S$ komutuje z $\partial$,
\[
 \begin{aligned}
 \partial\circ T^{(r)}+T^{(r)}\circ\partial
 &=\sum_{j=0}^{r-1}S^j\circ(\partial\circ T+T\circ\partial)\\
 &=\sum_{j=0}^{r-1}(S^{j+1}-S^j)=S^r-\id.
 \end{aligned}
\]

Średnice części po kolejnych podziałach
barycentrycznych dążą do zera: dla
$n$-sympleksu każdy nowy wierzchołek
jest średnią wierzchołków pewnej ściany,
a nowa krawędź łączy barycentra zagnieżdżonych ścian.
Jeśli mają one $p<q\leq n+1$ wierzchołków, to
$b_G=(p/q)b_F+((q-p)/q)b_{G\setminus F}$.
Zatem $\|b_G-b_F\|\leq(q-p)\operatorname{diam}(\Delta)/q
\leq n\operatorname{diam}(\Delta)/(n+1)$.
Średnica wypukłego sympleksu jest największą odległością jego
wierzchołków, więc dla $n\geq1$ otrzymujemy geometryczny spadek
średnic (także na ścianach). Dla $n=0$ sympleks jest już punktem. Dla danego $\sigma$
przeciwobrazy $U_\alpha$ tworzą
otwarte pokrycie zwartego $\Delta^n$;
jego liczba Lebesgue'a jest dodatnia.
Po dostatecznie wielu podziałach
każda część trafia przez $\sigma$
do jednego $U_\alpha$.
Każdy łańcuch jest skończoną sumą,
więc można wybrać wspólne $r$
dla wszystkich jego sympleksów.

Jeśli $z$ jest cyklem, $S^r(z)$
jest mały i różni się od $z$
o $\partial(T^{(r)}(z))$; to dowodzi
surjektywności na homologii.
Jeśli mały cykl $z$ jest brzegiem
$c$ w pełnym kompleksie, podzielmy
$c$ tak wiele razy, by $S^r(c)$
był mały. Łańcuch $T^{(r)}(z)$
też jest mały, ponieważ każda
jego część leży w obrazie
jednego pierwotnego, małego sympleksu.
Równości $\partial c=z$ i
$\partial(T^{(r)}(z))=S^r(z)-z$ dają
\[
 \partial\bigl(S^r(c)-T^{(r)}(z)\bigr)=z.
\]
Zatem $z$ jest brzegiem już
w podkompleksie małych łańcuchów:
otrzymujemy iniektywność.
\end{proof}

Stożkowanie wypełnia każdy cykl dodatniego stopnia w wypukłym
sympleksie, a w stopniu zero tylko cykl o sumie współczynników zero.
W konstrukcji $T$ stopień zero załatwiamy przez $t_0=0$,
bo $S_0=\id$. Lokalność $T$ jest istotna dla dowodu iniektywności.

\begin{przyklad}[Co dzieli operator $S$]
\label{prz:podzial-singularny-8}
Dla drogi $\sigma:[0,1]\to X$ dwa odwzorowania afiniczne podziału
mają wzory $\alpha(t)=t/2$ i $\beta(t)=1-t/2$.
Zatem $S_1(\sigma)=\sigma\circ\alpha-\sigma\circ\beta$.
Druga połowa jest parametryzowana od $1$ do $1/2$, a jej
współczynnik wynosi $-1$. Bez wprowadzania jakiejkolwiek
relacji między różnymi generatorami dostajemy
\[
 \partial\bigl(S_1(\sigma)\bigr)
 =\bigl([\sigma(1/2)]-[\sigma(0)]\bigr)
  -\bigl([\sigma(1/2)]-[\sigma(1)]\bigr)
 =\partial\sigma.
\]
W wymiarze $2$ łańcuch ma sześć składników: wybieramy wierzchołek,
zawierającą go krawędź i cały trójkąt. Ich barycentra tworzą
uporządkowane wierzchołki małego sympleksu.
Rysunek~\ref{fig:podzial-barycentryczny-8} pokazuje ten podział
w dziedzinie. Obraz $\sigma(\Delta^2)$ może mieć samoprzecięcia;
nie trzeba go ciąć jako podzbioru $X$.
\end{przyklad}

\begin{figure}[!htbp]
\centering
\begin{tikzpicture}[>=Stealth,line cap=round,line join=round]
 \coordinate (a) at (0,0); \coordinate (b) at (4,0);
 \coordinate (c) at (1.6,3); \coordinate (mab) at (2,0);
 \coordinate (mac) at (.8,1.5); \coordinate (mbc) at (2.8,1.5);
 \coordinate (o) at (1.8667,1);
 \fill[manifoldblue!10] (a)--(b)--(c)--cycle;
 \draw[manifoldblue!70!black,thick] (a)--(b)--(c)--cycle;
 \foreach \p in {a,b,c,mab,mac,mbc}{\draw[manifoldred] (o)--(\p);}
 \foreach \p in {mab,mac,mbc,o}{\fill[manifoldgold!85!black] (\p) circle(2.1pt);}
 \node[below left] at (a) {$e_0$}; \node[below right] at (b) {$e_1$};
 \node[above] at (c) {$e_2$}; \node[below] at (mab) {$b_{01}$};
 \node[left] at (mac) {$b_{02}$}; \node[right] at (mbc) {$b_{12}$};
 \node[above=2pt,font=\small] at (o) {$b$};
 \node[align=left,text width=5.2cm,anchor=west] at (4.7,1.5)
  {Wierzchołki jednej części:\\[3pt]
   $e_0,\ b_{01},\ b$.\\[6pt]
   Sześć części odpowiada sześciu permutacjom
   wierzchołków wyjściowego trójkąta.};
\end{tikzpicture}
\caption{Podział barycentryczny standardowego $2$-sympleksu.
Wspólne ściany małych części znoszą się w $\partial\circ S$,
a na zewnętrznych krawędziach pozostaje podział $S\circ\partial$.}
\label{fig:podzial-barycentryczny-8}
\end{figure}

\begin{twierdzenie}[Wycinanie]
\label{thm:wycinanie}
Jeżeli $Z\subseteq A\subseteq X$ oraz
$\overline Z\subseteq\operatorname{int}_X A$,
to inkluzja par
\[
 (X\setminus Z,A\setminus Z)\longrightarrow(X,A)
\]
indukuje izomorfizm na wszystkich
grupach homologii względnej.
\end{twierdzenie}
\begin{proof}
Oznaczmy $U=X\setminus\overline Z$
oraz $V=\operatorname{int}_X A$.
To otwarte pokrycie $X$;
po ograniczeniu do $X\setminus Z$
otrzymujemy pokrycie
$\{U,V\setminus Z\}$.
Lemat~\ref{lem:male-lancuchy} działa
także po wzięciu ilorazu przez
łańcuchy w $A$: izomorfizmy
dla $A$ i $X$, zestawione
z długimi ciągami par, dają
izomorfizm dla odpowiadających
kompleksów względnych
(Lemat~\ref{lem:pieciu}).

Policzmy teraz te małe kompleksy.
Ponieważ $V\subseteq A$,
\[
 C_\bullet^{\{U,V\}}(X)
   =C_\bullet(U)+C_\bullet(V),\qquad
 C_\bullet^{\{U,V\}}(A)
   =C_\bullet(U\cap A)+C_\bullet(V).
\]
Ich iloraz jest kanonicznie
$C_\bullet(U)/C_\bullet(U\cap A)$:
sympleks należący jednocześnie do
$U$ i $V$ należy do $U\cap A$,
więc jądro tej identyfikacji jest
dokładnie wskazanym podkompleksem.
Analogiczny iloraz dla
$(X\setminus Z,A\setminus Z)$
i jego pokrycia jest \emph{tym samym}
kompleksem
$C_\bullet(U)/C_\bullet(U\cap A)$.
Inkluzja par staje się na tych
modelach identycznością, co
kończy dowód.
\end{proof}

Wycinanie mówi, że usunięcie części
leżącej głęboko we wnętrzu $A$
nie zmienia tego, co para $(X,A)$
mierzy \emph{poza} $A$. Warunek
o domknięciu ma znaczenie:
zapewnia otwarte pokrycia użyte
w dowodzie.

Nie można dowolnie usuwać punktów $A$. Dla
$(X,A)=([0,1],\{0\})$ wszystkie grupy względne znikają,
jak w Przykładzie~\ref{prz:para-odcinek-konce}.
Jeśli wytniemy $Z=\{0\}$, otrzymamy parę $((0,1],\varnothing)$,
której $H_0$ jest równe $\Z$. Tutaj $\operatorname{int}_X A$
jest puste, więc hipoteza wycinania nie zachodzi.
Usunęliśmy cały wyróżniony zbiór, zamiast tylko jego wewnętrznej części.

\begin{twierdzenie}[Ciąg Mayera--Vietorisa]
\label{thm:mv-przestrzenie}
Jeśli $X=U\cup V$, gdzie $U,V$
są otwarte w $X$, to istnieje
długi ciąg dokładny
\begin{equation}\label{eq:mv-przestrzenie}
 \cdots\to H_n(U\cap V)
 \xrightarrow{(i_*,-j_*)}
 H_n(U)\oplus H_n(V)
 \xrightarrow{k_*+\ell_*}H_n(X)
 \xrightarrow{\delta}H_{n-1}(U\cap V)\to\cdots.
\end{equation}
Tak samo działają grupy zredukowane.
\end{twierdzenie}
\begin{proof}
Wewnątrz $C_\bullet(X)$ rozważmy
$B_\bullet=C_\bullet(U)+C_\bullet(V)$.
Mamy krótki ciąg kompleksów
\[
 0\to C_\bullet(U\cap V)
 \xrightarrow{c\mapsto(c,-c)}
 C_\bullet(U)\oplus C_\bullet(V)
 \xrightarrow{(a,b)\mapsto a+b}
 B_\bullet\to0.
\]
Iniektywność pierwszego odwzorowania jest jasna;
jądro drugiej składa się z par
$(c,-c)$, w których każdy singularny
sympleks występujący w $c$ należy
zarazem do $U$ i do $V$.
Surjektywność wynika z definicji sumy.
Ponadto $B_\bullet=C_\bullet^{\{U,V\}}(X)$,
więc Lemat~\ref{lem:male-lancuchy}
identyfikuje $H_*(B)$ z $H_*(X)$.
Zastosowanie Twierdzenia~\ref{thm:les-homologia}
daje~\eqref{eq:mv-przestrzenie}.
W wersji zredukowanej dołączamy
do każdego kompleksu stopień $-1$
równy $\Z$ i różniczkę
$\varepsilon:C_0\to\Z$.
Na stopniu $-1$ ten sam krótki
ciąg ma postać
$0\to\Z\xrightarrow{c\mapsto(c,-c)}
\Z\oplus\Z\xrightarrow{(a,b)\mapsto a+b}
\Z\to0$, nadal dokładną.
W stopniach nieujemnych
obowiązuje poprzedni rachunek.
Homologie tych powiększonych
kompleksów są właśnie
homologiami zredukowanymi;
dotyczy to również pustego
przecięcia, dla którego
$\widetilde H_{-1}(\varnothing)=\Z$.
\end{proof}

Jeśli $c=a+b\in B_n$ jest cyklem,
gdzie $a\in C_n(U)$ i $b\in C_n(V)$,
to $\partial a=-\partial b$ jest
łańcuchem w $U\cap V$, a
$\delta[c]=[\partial a]$ z wyborem
znaku zgodnym z pierwszym odwzorowaniem.
To praktyczna recepta na odwzorowanie łączące.
Istotnie, parę $(a,b)$ podnoszącą $c$ różniczka przenosi
na $(\partial a,\partial b)=(\partial a,-\partial a)$,
czyli obraz $\partial a$ przez pierwszą strzałkę krótkiego ciągu.
Jeżeli $c$ nie jest jeszcze sumą takich małych sympleksów,
najpierw zastępujemy go homologicznie przez $S^r(c)$.

\begin{przyklad}[Okrąg z dwóch łuków]
\label{prz:mv-okrag}
Pokryjmy $S^1$ dwoma otwartymi,
ściągalnymi łukami $U,V$, których
przecięcie składa się z dwóch
łuków. Mamy
$\widetilde H_0(U)=
\widetilde H_0(V)=0$,
a $\widetilde H_0(U\cap V)\cong\Z$
generuje różnica klas punktów
z dwóch składowych.
Zredukowany ciąg
Mayera--Vietorisa daje
\[
 0=\widetilde H_1(U)\oplus
 \widetilde H_1(V)\to
 \widetilde H_1(S^1)
 \xrightarrow[\cong]{\delta}
 \widetilde H_0(U\cap V)\to0.
\]
Zatem $H_1(S^1)\cong\Z$.
Geometria odwzorowania $\delta$ jest prosta:
rozcinając cykl okręgu na łuki
leżące w $U$ i $V$, otrzymujemy
dwa końce w różnych składowych
ich przecięcia.
Wybierzmy punkty $p,q$ w tych składowych, drogę $a$ z $p$ do $q$
w $U$ i drogę $b$ z $q$ do $p$ w $V$.
Łańcuch $c=a+b$ jest cyklem, a
$\delta[c]=[q]-[p]$ generuje $\widetilde H_0(U\cap V)$.
Zatem $[c]$ jest generatorem $H_1(S^1)$, a nie jedynie
nieokreśloną wielokrotnością generatora.
\end{przyklad}

\begin{figure}[!htbp]
\centering
\begin{tikzpicture}[>=Stealth,line cap=round,line join=round]
 \draw[manifoldblue!15,line width=13pt] (20:1.65) arc (20:340:1.65);
 \draw[manifoldgreen!17,line width=7pt] (200:1.65) arc (200:520:1.65);
 \draw[manifoldblue!80!black,thick,->] (90:1.65) arc (90:270:1.65);
 \draw[manifoldgreen!70!black,thick,->] (270:1.65) arc (270:450:1.65);
 \fill[manifoldred] (90:1.65) circle(2.8pt) node[above=5pt] {$p$};
 \fill[manifoldred] (270:1.65) circle(2.8pt) node[below=5pt] {$q$};
 \node[manifoldblue!80!black,left] at (-1.7,0) {$a\subset U$};
 \node[manifoldgreen!70!black,right] at (1.7,0) {$b\subset V$};
 \node[align=left,text width=4.8cm,anchor=west] at (3.15,.15)
  {$U\cap V$ ma dwie składowe:\\[3pt]
   jedną wokół $p$, drugą wokół $q$.\\[7pt]
   $\partial a=[q]-[p]$,\\[3pt]
   $\partial b=[p]-[q]$.};
\end{tikzpicture}
\caption{Pokrycie okręgu dwoma otwartymi łukami. Niebieski łuk
omija punkt po prawej, zielony omija punkt po lewej.
Grubsze pasy zaznaczają zbiory pokrycia; strzałki wskazują
dwie drogi tworzące cykl.}
\label{fig:mv-okrag-8}
\end{figure}

\begin{twierdzenie}[Homologia sfer]
\label{thm:homologia-sfer}
Dla $n\geq0$:
\[
 \widetilde H_k(S^n;\Z)=
 \begin{cases}\Z,&k=n,\\0,&k\ne n.\end{cases}
\]
Ponadto dla $n\geq1$
\[
 H_k(D^n,S^{n-1};\Z)=
 \begin{cases}\Z,&k=n,\\0,&k\ne n.\end{cases}
\]
\end{twierdzenie}
\begin{proof}
$S^0$ ma dwie składowe punktowe:
$H_0(S^0)=\Z^2$ oraz
$\widetilde H_0(S^0)=\Z$.
Wyższa homologia znika, bo jest
to suma dwóch kopii homologii
punktu z Przykładu~\ref{prz:homologia-punktu}.

Dla $n\geq1$ weźmy lekko
powiększone otwarte półkule
$U,V\subset S^n$. Konkretnie, dla $0<\varepsilon<1$ i ostatniej
współrzędnej $z$ możemy wybrać $U=\{z>-\varepsilon\}$
i $V=\{z<\varepsilon\}$; przecięcie to pas $|z|<\varepsilon$.
Każda półkula jest
ściągalna, a $U\cap V$
deformacyjnie zwija się na
sferę równikową $S^{n-1}$.
Można użyć homotopii wysokości z
Przykładu~\ref{prz:pi1-sfera}: w pasie przecięcia
zmiana $z$ na $(1-t)z$ pozostawia punkt w tym pasie,
a dla $t=1$ daje równik.
Zredukowany ciąg
Mayera--Vietorisa, wraz z
zerowaniem homologii $U,V$,
daje dla każdego $k$ izomorfizm
\[
 \widetilde H_k(S^n)
 \xrightarrow{\ \cong\ }
 \widetilde H_{k-1}(S^{n-1}).
\]
W stopniu $k=1$ używa się
właśnie zredukowanego $H_0$;
gdy $n=1$, obie składowe
przecięcia są niezbędne.
Indukcja od $S^0$ daje pierwszy wzór.

Dysk $D^n$ jest ściągalny. W ciągu
zredukowanym pary $(D^n,S^{n-1})$
grupy dysku są zerowe, a odwzorowanie
łączące daje
$H_k(D^n,S^{n-1})\cong
\widetilde H_{k-1}(S^{n-1})$
dla $k\geq1$. Dla $k=0$
odwzorowanie $H_0(S^{n-1})\to H_0(D^n)$
jest surjekcją, więc grupa
względna również znika.
To daje drugi wzór.
\end{proof}

\begin{uwaga}[Klasa fundamentalna i orientacja]
Generator $H_n(D^n,S^{n-1})\cong\Z$
nazywamy \emph{klasą orientacji dysku}.
Wybór jednego z jego dwóch znaków
określa orientację; w długim ciągu
pary $\delta$ posyła ją na
generator $\widetilde H_{n-1}(S^{n-1})$.
Różniczkowe ujęcie orientacji
i całkowania pozostawiamy
osobnemu rozdziałowi.
\end{uwaga}

\section{Od triangulacji do homologii przestrzeni}
\label{sec:porownanie-symplicjalne}

Dla kompleksu symplicjalnego zachowujemy ustalony porządek
wierzchołków, a dla $\Delta$-kompleksu porządki jego odwzorowań charakterystycznych.
Każdy bazowy sympleks $\sigma:\Delta^n\to|K|$ jest wtedy
jednocześnie singularnym sympleksem. Daje to odwzorowanie łańcuchowe
\[
 I:C_\bullet^\Delta(K)\longrightarrow C_\bullet(|K|),
 \qquad [\sigma]\longmapsto \sigma,
\]
bo oba operatory $\partial$ sumują te same
ściany z tymi samymi znakami.
Odwzorowanie $I$ rozszerzamy liniowo: na przykład $I(-[\sigma])=-\sigma$.
Nie zastępujemy prawej strony osobnym generatorem otrzymanym
przez odwrócenie parametryzacji. Nie jest to
izomorfizm grup łańcuchów: po prawej są
wszystkie odwzorowania ciągłe, po lewej tylko
wybrane sympleksy. Niezwykłe jest to,
że po przejściu do homologii różnica znika.

Potrzebny będzie prosty sposób traktowania
komórki z całym brzegiem zwiniętym do punktu.

\begin{lemat}[Iloraz pary komórkowej]
\label{lem:iloraz-cw-homologia}
Niech $A\ne\varnothing$ będzie podkompleksem skończonego
kompleksu CW $X$. Jeśli istnieje otwarte
otoczenie $N$ zbioru $A$, które silnie
deformacyjnie retrahuje się na $A$, to
odwzorowanie ilorazowe indukuje izomorfizm
\[
 H_k(X,A)\xrightarrow{\ \cong\ }
 H_k(X/A,\{*\})=\widetilde H_k(X/A)
 \quad(k\geq0).
\]
Takie otoczenie istnieje, gdy
$A=X^{n-1}\ne\varnothing$ i $X=X^n$ jest skończonym
szkieletem CW dla $n\geq1$, a także gdy do skończonego $A$
dołączamy pojedynczą komórkę dodatniego wymiaru.
\end{lemat}
\begin{proof}
Inkluzja $A\hookrightarrow N$ jest
równoważnością homotopijną. Porównując
długie ciągi par $(X,A)$ i $(X,N)$
za pomocą Lematu~\ref{lem:pieciu},
dostajemy $H_k(X,A)\cong H_k(X,N)$.
Po zwinięciu $A$ do punktu homotopia schodzi na $N/A$.
Uzasadnijmy jej ciągłość również względem czasu.
Skończony kompleks CW jest zwarty, bo jest obrazem skończonej
sumy zwartych dysków; domknięty $A$ też jest zwarty.
Projekcja $q:N\to N/A$ jest domknięta: dla domkniętego $F\subseteq N$
zbiór $q^{-1}q(F)$ jest równy $F$ lub $F\cup A$.
Ma zwarte włókna, więc $q\times\id_I$ jest również domknięta.
Istotnie, jeśli domknięty zbiór omija włókno nad $(y,t)$,
skończone pokrycie tego włókna prostokątami otwartymi daje
otoczenia $O\supseteq q^{-1}(y)$ i $J\ni t$ z $O\times J$
rozłącznym z tym zbiorem; domkniętość $q$ daje otoczenie
$W\ni y$ takie, że $q^{-1}(W)\subseteq O$.
Stąd $(q\times\id_I)$ jest odwzorowaniem ilorazowym,
co pozwala przenieść homotopię. Zatem $N/A$ silnie zwija się do $\{*\}$,
więc tak samo
$H_k(X/A,\{*\})\cong H_k(X/A,N/A)$.

Porównajmy środkowe grupy.
Zbiór $A$ jest domknięty
(dla skończonego CW wynika to
z topologii słabej), zatem
$\{X\setminus A,N\}$ jest
otwartym pokryciem $X$.
W ilorazie mamy pokrycie
$\{(X/A)\setminus\{*\},N/A\}$.
Lemat o małych łańcuchach,
zastosowany również do $N$,
pozwala obliczać homologie
tych dwóch par odpowiednio za pomocą kompleksów
\[
 \frac{C_\bullet(X\setminus A)}
      {C_\bullet(N\setminus A)}
 \quad\text{oraz}\quad
 \frac{C_\bullet((X/A)\setminus\{*\})}
      {C_\bullet((N/A)\setminus\{*\})}.
\]
Poza $A$ odwzorowanie ilorazowe jest
homeomorfizmem, więc te dwa
kompleksy są izomorficzne.
Składając trzy izomorfizmy,
otrzymujemy tezę.

Dla $A=X^{n-1}\subset X^n$
weźmy $N$ złożony z $A$ oraz,
w każdej dołączonej $n$-komórce,
obrazu pierścienia
$\{x\in D^n:\|x\|>1/2\}$.
Jest otwarty w topologii
nadanej przez odwzorowania dołączające.
W każdym pierścieniu promieniowo
przesuwamy $r$ od wartości
$r>1/2$ do $1$, pozostawiając
brzeg i cały $A$ nieruchome.
Na częściach dyskowych używamy wzoru
$ru\mapsto((1-t)r+t)u$. Ciągłość przy $A$ sprawdzamy jak
w dowodzie Twierdzenia~\ref{thm:pi1-dolaczenie-komorki}:
w każdym otoczeniu zachowujemy tylko te punkty pasa, których
cały odcinek radialny do brzegu mieści się w tym otoczeniu.
Te mniejsze zbiory są otwarte i pozostają w pierwotnym otoczeniu
przez cały czas homotopii. Daje to silną retrakcję $N\to A$.
Ten sam argument działa przy dołączeniu pojedynczej komórki.
\end{proof}

\begin{lemat}[Generator pojedynczego sympleksu]
\label{lem:orientacja-simpleksu}
Identycznościowe odwzorowanie singularne
$u_n:\Delta^n\to\Delta^n$ określa
generator
$H_n(\Delta^n,\partial\Delta^n)\cong\Z$
dla $n\geq1$.
\end{lemat}
\begin{proof}
Jest cyklem względnym, bo jego
brzeg leży w $\partial\Delta^n$.
Brzeg sympleksu jest sferą
(Stwierdzenie~\ref{prop:brzeg-sympleksu}),
a jej homologię policzyliśmy
w Twierdzeniu~\ref{thm:homologia-sfer}.
Wynika stąd, że odwzorowanie łączące
\[
 \delta:H_n(\Delta^n,\partial\Delta^n)
 \longrightarrow
 \widetilde H_{n-1}(\partial\Delta^n)
 \]
jest izomorfizmem i posyła
$[u_n]$ na klasę cyklu
$\partial u_n$.
Pokażemy indukcyjnie, że ten
cykl jest generatorem.
Dla $n=1$ jest to różnica
końców, generator
$\widetilde H_0(S^0)$.
W wyższym wymiarze rozdzielmy
$\partial\Delta^n$ na jedną
ścianę $F\cong\Delta^{n-1}$
i sumę $B$ pozostałych ścian.
$F$ i $B$ są ściągalne:
$B$ jest stożkiem z przeciwległego
wierzchołka nad $\partial F$;
ich przecięciem jest $\partial F$.
W ciągu Mayera--Vietorisa
odwzorowanie łączące
$\widetilde H_{n-1}(\partial\Delta^n)
\to\widetilde H_{n-2}(\partial F)$
jest izomorfizmem. Aby zastosować wersję dla zbiorów otwartych, przedstawmy
oba dyski $F,B$ jako stożki nad ich wspólnym brzegiem $\partial F$.
Niech $r\in[0,1]$ oznacza w każdym z nich parametr radialny,
z $r=1$ na $\partial F$. Zbiory
$U=F\cup\{x\in B:r(x)>1/2\}$ oraz
$V=B\cup\{x\in F:r(x)>1/2\}$ są otwarte w $F\cup B$:
na każdej z dwóch części ich przeciwobrazy są otwarte.
Przesuwając dodatkowe pasy do $r=1$, retrahujemy je na $F,B$,
a ich przecięcie na $\partial F$.
Cykl $\partial u_n$ jest sumą
$\pm u_F$ na $F$ i pozostałych
ścian na $B$; jego obrazem przez
odwzorowanie łączące jest
$\pm[\partial u_F]$.
To generator na mocy
założenia indukcyjnego.
Zatem $\partial u_n$, a więc
również $u_n$ względnie,
jest generatorem.
\end{proof}

\begin{twierdzenie}[Porównanie homologii]
\label{thm:symplicjalna-singularna}
Dla skończonego kompleksu symplicjalnego
$K$ naturalne odwzorowanie indukuje
\[
 I_*:H_n^\Delta(K)\xrightarrow{\ \cong\ }
 H_n(|K|;\Z)
 \qquad(n\geq0).
\]
To samo zachodzi dla skończonego
$\Delta$-kompleksu, jeśli generatory
po lewej to jego sympleksy
charakterystyczne z ich ścianami.
\end{twierdzenie}
\begin{proof}
Uporządkujmy sympleksy tak, aby
każdy pojawiał się po wszystkich
swoich ścianach. Budujemy kompleks
po jednym sympleksie:
$L\subset K=L\cup_\varphi\Delta^n$.
Załóżmy indukcyjnie, że porównanie
jest izomorfizmem dla $L$.
Iloraz kompleksów symplicjalnych
$C^\Delta(K)/C^\Delta(L)$
ma jeden generator w stopniu $n$
i zero w pozostałych stopniach,
więc jego homologia jest $\Z$
tylko w stopniu $n$.

Jeśli $n\geq1$, po zwinięciu $L$
do punktu przestrzeń $|K|/|L|$
jest $\Delta^n/\partial\Delta^n
\cong S^n$; analogiczny opis
działa dla $\Delta$-kompleksu,
nawet gdy jego ściany były
między sobą utożsamione.
Lemat~\ref{lem:iloraz-cw-homologia}
i Twierdzenie~\ref{thm:homologia-sfer}
dają $H_q(|K|,|L|)=\Z$ dla
$q=n$ i zero w innych stopniach.
Odwzorowanie względne przenosi jedyny
generator na klasę charakterystycznego sympleksu.
Odwzorowanie par $(\Delta^n,\partial\Delta^n)\to(|K|,|L|)$
indukuje homeomorfizm ich ilorazów: jest ciągłą bijekcją
ze zwartej przestrzeni do przestrzeni Hausdorffa.
Naturalność odwzorowań do ilorazów w Lemacie~\ref{lem:iloraz-cw-homologia}
pokazuje, że także na homologii względnej jest izomorfizmem.
Na mocy Lematu~\ref{lem:orientacja-simpleksu} przenosi więc
generator na generator, a nie na jego dowolną wielokrotność.
Dla $n=0$ nowy wierzchołek jest
osobną składową i odpowiednie
odwzorowanie względne jest izomorfizmem
$\Z\to\Z$ w stopniu $0$.

Mamy teraz przemienny diagram
długich ciągów par dla kompleksów
symplicjalnych i singularnych.
Oznaczając homologię ilorazu kompleksów symplicjalnych przez
$H_*^\Delta(K,L)$, dla dowolnego $q\geq0$ jego potrzebny fragment zapisujemy jako
\[
\begin{tikzcd}[column sep=small,row sep=large]
 H_{q+1}^\Delta(K,L) \arrow[r] \arrow[d,"\cong" left]
 &H_q^\Delta(L) \arrow[r] \arrow[d,"\cong" left]
 &H_q^\Delta(K) \arrow[r] \arrow[d,"I_*"]
 &H_q^\Delta(K,L) \arrow[r] \arrow[d,"\cong" right]
 &H_{q-1}^\Delta(L) \arrow[d,"\cong" right]\\
 H_{q+1}(|K|,|L|) \arrow[r]
 &H_q(|L|) \arrow[r]
 &H_q(|K|) \arrow[r]
 &H_q(|K|,|L|) \arrow[r]
 &H_{q-1}(|L|).
\end{tikzcd}
\]
Przemienność przy odwzorowaniu łączącym wynika z tego, że
$I$ posyła brzeg generatora dokładnie na jego brzeg singularny.
Na wyrazach dotyczących $L$
odwzorowania pionowe są izomorfizmami
z założenia indukcyjnego,
a na względnych właśnie
to pokazaliśmy. Lemat~\ref{lem:pieciu}
zastosowany do kolejnych pięciu
wyrazów daje izomorfizm dla $K$.
Zaczynamy od kompleksu pustego,
dla którego wszystkie zwykłe
grupy homologii są zerowe.
\end{proof}

W praktyce wolno zatem obliczać
homologię przestrzeni z jednej
wygodnej triangulacji: otrzymany
wynik nie zależy od wyboru
triangulacji. Podziały sympleksów
zmieniają grupy łańcuchów,
ale nie grupy homologii.

Na przykład dla kwadratu z
Przykładu~\ref{prz:homologia-kwadrat} obliczone grupy
$H_0^\Delta=\Z$, $H_k^\Delta=0$ dla $k>0$ są już grupami
homologii całego wypełnionego kwadratu. To samo otrzymalibyśmy
z dowolnego drobniejszego podziału, mimo większych macierzy brzegu.
Nie twierdzimy przy tym, że każda przestrzeń ma triangulację:
homologia singularna pozostaje definicją także bez niej.

\section{Homologia komórkowa}
\label{sec:homologia-komorkowa}

W triangulacji torusa potrzeba
trójkątów i ich ścian. Struktura
CW z Sekcji~\ref{sec:cw-kompleksy}
może opisywać go jednym punktem,
dwoma pętlami i jedną
$2$-komórką. Homologia komórkowa
zachowuje tę oszczędność.

\begin{twierdzenie}[Względna homologia warstwy komórek]
\label{thm:warstwa-komorek}
Niech $X$ będzie skończonym
kompleksem CW, $X^n$ jego
$n$-szkieletem i $X^{-1}=\varnothing$.
Dla $n\geq1$:
\[
 H_k(X^n,X^{n-1})
 \cong
 \begin{cases}
 \displaystyle\bigoplus_{\alpha\in I_n}\Z[e_\alpha^n],
 &k=n,\\[3pt]0,&k\ne n,
 \end{cases}
\]
gdzie $I_n$ oznacza zbiór $n$-komórek.
Dla $n=0$ mamy
$H_0(X^0)=\bigoplus_{\alpha\in I_0}\Z[e_\alpha^0]$
i $H_k(X^0)=0$ dla $k\ne0$.
\end{twierdzenie}
\begin{proof}
Jeśli $X^{n-1}=\varnothing$ i $n\geq1$, to również $X^n=\varnothing$:
nie można dołączyć dysku o niepustym brzegu do pustej przestrzeni.
Teza jest wtedy oczywista. W pozostałych przypadkach
w ilorazie $X^n/X^{n-1}$
brzeg każdej dołączonej kopii
$D^n$ staje się jednym punktem,
a wnętrza komórek pozostają
rozłączne. Z definicji topologii
CW otrzymujemy bukiet
$\bigvee_{\alpha\in I_n}
(D^n/\partial D^n)\cong
\bigvee_{\alpha\in I_n}S^n$.
Lemat~\ref{lem:iloraz-cw-homologia}
zamienia homologię względną
na zredukowaną homologię bukietu.
Dla skończonego bukietu
zredukowana homologia jest sumą
zredukowanych homologii składników:
dołączamy kolejne sfery. Dla bukietu $Y\vee S^n$
bierzemy $U=Y\cup B$ i $V=S^n\cup N$, gdzie $B$ jest małą
otwartą kulą wokół punktu wspólnego w nowej sferze, a $N$
jest sumą takich kul w dotychczasowych sferach.
Te zbiory są otwarte w topologii ilorazowej bukietu,
retrahują się na $Y,S^n$, a $U\cap V=N\vee B$ jest ściągalne.
Ciąg Mayera--Vietorisa daje izomorfizm indukowany inkluzjami
składników, a więc właśnie ich sumę prostą.
Twierdzenie~\ref{thm:homologia-sfer}
daje wskazane grupy.
Znaki generatorów zależą
od wyboru orientacji komórek.
Dla $n=0$ mamy skończony
zbiór punktów; rachunek wynika
z homologii punktu.
\end{proof}

\begin{definicja}[Kompleks komórkowy]
\index{kompleks komorkowy@Kompleks komórkowy}
\label{def:kompleks-komorkowy}
Kładziemy
$C_n^{\mathrm{CW}}(X)=H_n(X^n,X^{n-1})$.
Dla $n\geq1$ jego operator brzegu
jest złożeniem
\begin{equation}\label{eq:cellular-boundary}
 H_n(X^n,X^{n-1})
 \xrightarrow{\ \delta_n\ }
 H_{n-1}(X^{n-1})
 \xrightarrow{\ j_{n-1}\ }
 H_{n-1}(X^{n-1},X^{n-2}),
\end{equation}
gdzie $\delta_n$ pochodzi z
długiego ciągu pary, a
$j_{n-1}$ z ilorazu łańcuchów.
Przyjmujemy $\partial_0^{\mathrm{CW}}=0$.
\end{definicja}

Pierwsza strzałka bierze brzeg komórki jako klasę w całym
poprzednim szkielecie. Druga zapomina jego część w $X^{n-2}$,
zostawiając całkowite krotności przejść przez $(n-1)$-komórki.
Zatem $\partial_n^{\mathrm{CW}}=j_{n-1}\circ\delta_n$;
nie trzeba wybierać triangulacji wnętrza każdej komórki.

\begin{twierdzenie}[Komórki obliczają homologię]
\label{thm:komorkowa-singularna}
Operator~\eqref{eq:cellular-boundary}
spełnia $(\partial^{\mathrm{CW}})^2=0$.
Dla skończonego CW istnieją
izomorfizmy
\[
 H_n(C_\bullet^{\mathrm{CW}}(X),
       \partial^{\mathrm{CW}})
 \cong H_n(X;\Z)
 \qquad\text{dla każdego }n\geq0.
\]
\end{twierdzenie}
\begin{proof}
Dla $n=1$ równość brzegu brzegu wynika z $\partial_0^{\mathrm{CW}}=0$.
Dla $n\geq2$ z dokładności długiego ciągu
pary $(X^{n-1},X^{n-2})$
wynika $\delta_{n-1}\circ j_{n-1}=0$.
Stąd
\[
 \partial^{\mathrm{CW}}_{n-1}\circ
 \partial^{\mathrm{CW}}_n
 =j_{n-2}\circ\delta_{n-1}\circ j_{n-1}\circ\delta_n=0.
\]

Rozpiszmy porównanie homologii.
Indukcja po szkieletach,
z użyciem długiego ciągu pary
i Twierdzenia~\ref{thm:warstwa-komorek},
daje $H_k(X^m)=0$ dla $k>m$.
Inkluzja $X^{m-1}\hookrightarrow X^m$
indukuje ponadto izomorfizm na
$H_k$ dla $k<m-1$: dwa sąsiednie
wyrazy względne o stopniach
$k+1$ i $k$ są wtedy zerowe.
W szczególności komórki wymiaru
$\geq n+2$ nie zmieniają $H_n$.

Dla $n\geq1$ w ciągu pary
$(X^n,X^{n-1})$ mamy odcinek
\[
 0=H_n(X^{n-1})\longrightarrow
 H_n(X^n)\xrightarrow{\ j_n\ }
 C_n^{\mathrm{CW}}
 \xrightarrow{\ \delta_n\ }
 H_{n-1}(X^{n-1}).
\]
Odwzorowanie $j_{n-1}:H_{n-1}(X^{n-1})
\to C_{n-1}^{\mathrm{CW}}$
jest iniektywne: jego jądro
jest obrazem
$H_{n-1}(X^{n-2})=0$.
Dla $n=1$ używamy
$X^{-1}=\varnothing$, co daje
ten sam wniosek.
Zatem $\ker\partial_n^{\mathrm{CW}}
=\ker\delta_n=\operatorname{im}j_n$
i $j_n$ identyfikuje cykle
komórkowe z $H_n(X^n)$.

W następnym długim ciągu,
dla $(X^{n+1},X^n)$,
\[
 C_{n+1}^{\mathrm{CW}}
 \xrightarrow{\ \delta_{n+1}\ }
 H_n(X^n)\longrightarrow
 H_n(X^{n+1})\longrightarrow
 H_n(X^{n+1},X^n)=0.
\]
Jądro przejścia do $H_n(X^{n+1})$
jest więc $\operatorname{im}\delta_{n+1}$.
Po włożeniu w $C_n^{\mathrm{CW}}$
przez $j_n$ jest to dokładnie
$\operatorname{im}\partial_{n+1}^{\mathrm{CW}}$.
Otrzymujemy
$H_n^{\mathrm{CW}}(X)\cong
H_n(X^{n+1})\cong H_n(X)$.
Dla $n=0$ ten sam argument działa
z $H_0(X^0)=C_0^{\mathrm{CW}}$:
$H_0(X^1)$ jest ilorazem przez
obrazy brzegów $1$-komórek,
a wyższe komórki już nie
zmieniają $H_0$.
\end{proof}

W dowodzie pojawiają się trzy kolejne szkielety:
$X^{n-1}\subseteq X^n\subseteq X^{n+1}$.
Najpierw ustalamy, które kombinacje $n$-komórek mają zerowy brzeg;
następnie dzielimy je przez brzegi $(n+1)$-komórek.
Komórki jeszcze wyższych wymiarów nie zmieniają już $H_n$.
Nie oznacza to, że każda dołączona komórka daje nowy generator
homologii: może ona również usunąć wcześniejszą klasę.

\begin{przyklad}[Wypełnienie pętli i powstanie sfery]
\label{prz:dwa-dyski-homologia-8}
Zacznijmy od jednego okręgu z komórkami $v,a$.
Dołączmy dysk $e$ przez identyczność jego brzegu z okręgiem.
Wybieramy orientację tak, by $\partial_2e=a$.
Otrzymana przestrzeń jest dyskiem, a mnożenie przez $1$
z $C_2=\Z e$ do $C_1=\Z a$ zabija klasę $[a]$.

Do tej samej pętli dołączmy drugi dysk $f$, również z
$\partial_2f=a$. Dwa dyski sklejone całymi brzegami tworzą $S^2$.
Teraz $\partial_2(pe+qf)=(p+q)a$, więc
$\ker\partial_2=\Z(e-f)$ i $\operatorname{im}\partial_2=\Z a$.
Nie ma $3$-komórek: $H_2=\Z[e-f]$, podczas gdy $H_1=0$.
Różnica $e-f$ oznacza zgodne złożenie orientacji obu półsfer:
ich brzegi się znoszą. To samo dołączenie $2$-komórki może więc
usunąć klasę w wymiarze $1$ albo utworzyć klasę w wymiarze $2$,
zależnie od tego, czy jej brzeg był już brzegiem wcześniejszego łańcucha.
\end{przyklad}

\section{Odwzorowania dołączające i konkretne obliczenia}
\label{sec:obliczenia-homologii}

Znamy już grupy łańcuchów komórkowych: ich bazy stanowią komórki.
Do obliczenia homologii brakuje jeszcze macierzy brzegu.
Odczytamy jej współczynniki z odwzorowań dołączających,
a następnie oddzielnie wyznaczymy jądra i obrazy.
W ten sposób sklejenie przestrzeni staje się
konkretnym rachunkiem nad $\Z$.

\begin{definicja}[Stopień odwzorowania sfer]
\index{stopien odwzorowania sfer@Stopień odwzorowania sfer}
\label{def:stopien-sfer}
Dla $m\geq1$ każde ciągłe odwzorowanie
$h:S^m\to S^m$ indukuje
endomorfizm $H_m(S^m)\cong\Z$.
Istnieje więc dokładnie jedna
liczba całkowita $\deg h$ taka,
że $h_*([S^m])=(\deg h)[S^m]$.
Nazywamy ją \emph{stopniem} $h$.
\end{definicja}

Wybieramy ten sam generator $[S^m]$ w dziedzinie i przeciwdziedzinie.
Zastąpienie go przez $-[S^m]$ po obu stronach nie zmienia stopnia.
Liczba $\deg h$ mierzy algebraiczną krotność pokrycia sfery:
przejścia zachowujące orientację i odwracające ją wnoszą przeciwne znaki.
Nie jest to na ogół liczba punktów dowolnego włókna.

Z definicji i funktorialności
$\deg(\id)=1$ oraz
$\deg(h\circ g)=\deg h\deg g$.
Homotopijne odwzorowania mają ten sam
stopień. Odwzorowanie stałe ma stopień zero, ponieważ
na dodatniej homologii przechodzi przez grupę homologii punktu,
która jest zerowa. Na $S^1=\{z\in\C:|z|=1\}$
odwzorowanie $z\mapsto z^r$ ma stopień $r\in\Z$.
Uzasadnijmy przejście od liczby obiegów do homologii.
Pętla $\gamma(t)=e^{2\pi it}$ reprezentuje generator $H_1(S^1)$:
po podziale na łuki jej obraz przez izomorfizm Mayera--Vietorisa
z Przykładu~\ref{prz:mv-okrag} jest różnicą punktów dwóch składowych przecięcia.
Dla dowolnych dwóch pętli $\alpha,\beta$ opartych w tym samym punkcie
mamy $[\alpha*\beta]=[\alpha]+[\beta]$ w $H_1$.
Istotnie, złożenie $\alpha*\beta$ z afinicznym odwzorowaniem
$\Delta^2\to[0,1]$ o wartościach $0,1/2,1$ w kolejnych wierzchołkach
jest singularnym $2$-sympleksem o brzegu
$\beta-(\alpha*\beta)+\alpha$.
Ponadto homotopia pętli względem końców daje homologiczne cykle
przez wzór pryzmatyczny; pętla stała jest brzegiem stałego $2$-sympleksu.
Stąd $[\bar\alpha]=-[\alpha]$, bo $\alpha*\bar\alpha$ jest
homotopijna do stałej. Pętla $t\mapsto e^{2\pi irt}$ jest,
z dokładnością do zmiany parametryzacji, złożeniem $r$ obiegów dla $r>0$,
$|r|$ obiegów przeciwnych dla $r<0$
i jest stała dla $r=0$. Jej klasa to zatem $r[\gamma]$.
W szczególności $z\mapsto z^3$ ma stopień $3$,
a $z\mapsto\bar z=z^{-1}$ stopień $-1$.

\begin{twierdzenie}[Współczynniki brzegu komórkowego]
\label{thm:formula-komorkowa}
Niech $X$ będzie skończonym kompleksem CW, a
$\varphi_\alpha:S^{n-1}\to X^{n-1}$
będzie odwzorowaniem dołączającym
zorientowanej $n$-komórki
$e_\alpha^n$, gdzie $n\geq2$.
Dla każdej zorientowanej
$(n-1)$-komórki $e_\beta^{n-1}$
zwińmy w $X^{n-1}$ do punktu
dopełnienie jej otwartego wnętrza.
Otrzymamy odwzorowanie
\[
 q_\beta:X^{n-1}\longrightarrow
 X^{n-1}/(X^{n-1}\setminus e_\beta^{n-1})
 \cong S^{n-1}.
\]
Sferę docelową orientujemy zgodnie z komórką $e_\beta^{n-1}$,
a dziedzinę $\varphi_\alpha$ jako brzeg zorientowanej komórki $e_\alpha^n$.
Wówczas
\begin{equation}\label{eq:formula-komorkowa}
 \partial_n^{\mathrm{CW}}(e_\alpha^n)
 =\sum_{\beta}
 \deg(q_\beta\circ\varphi_\alpha)
 e_\beta^{n-1}.
\end{equation}
W stopniu $1$ brzeg zorientowanej
komórki od początku $v_-$ do końca
$v_+$ wynosi $\partial_1e=v_+-v_-$.
\end{twierdzenie}
\begin{proof}
Generator $e_\alpha^n$ jest
obrazem klasy orientacji
$(D^n,S^{n-1})$ w parze
$(X^n,X^{n-1})$. Pierwsza
strzałka~\eqref{eq:cellular-boundary}
zwraca klasę obrazu
zorientowanego brzegu $S^{n-1}$
przez $\varphi_\alpha$:
tak działa odwzorowanie łączące,
bo na łańcuchach bierze
dosłowny brzeg podniesionego
cyklu względnego.
Druga strzałka przenosi tę
klasę do
$H_{n-1}(X^{n-1},X^{n-2})$.
Według
Twierdzenia~\ref{thm:warstwa-komorek}
jest to suma kopii $\Z$,
po jednej dla każdego $\beta$.
Rzutowanie na składnik $\beta$
jest, poprzez iloraz z
Lematu~\ref{lem:iloraz-cw-homologia},
tym samym co $q_{\beta*}$.
Zatem współczynnik przy
$e_\beta^{n-1}$ to liczba,
przez którą
$(q_\beta\circ\varphi_\alpha)_*$
mnoży generator homologii
$S^{n-1}$: z definicji
jest to jego stopień.
Dla $n=1$ używamy
bezpośrednio
$\partial[0,1]=[1]-[0]$;
oba końce mogą w $X^0$
oznaczać ten sam wierzchołek.
\end{proof}

Warto przeczytać wzór geometrycznie:
liczymy, ile razy, z jakimi znakami,
brzeg wyższej komórki przechodzi
przez daną komórkę niższą. Aby policzyć współczynnik przy
$e_\beta^{n-1}$, pomijamy pozostałe komórki przez złożenie
$q_\beta\circ\varphi_\alpha$. Wynikiem jest pojedyncza liczba całkowita,
czyli jeden wpis macierzy $\partial_n^{\mathrm{CW}}$.
Brzeg dołączanej komórki staje się brzegiem łańcucha w nowej przestrzeni,
więc jego klasa homologii
znika. Nie musi przy tym zniknąć każdy składnik tego brzegu:
relacja $2a=0$ pozostawia możliwość niezerowej klasy $[a]$ rzędu dwa.
Rysunek~\ref{fig:sklejenia-homologia-torus-rp2} zestawia
dwa sposoby dołączenia dysku, które dają różne wyniki.

\begin{figure}[!htbp]
\centering
\begin{tikzpicture}[>=Stealth,line cap=round,line join=round,
 every node/.style={font=\small}]
 % Torus: poziome boki mają kierunek w prawo, pionowe w górę.
 \node[font=\bfseries] at (1.7,4.1) {Torus $T^2$};
 \fill[manifoldblue!10] (0,0) rectangle (3.4,3.4);
 \draw[manifoldblue,very thick] (0,0) rectangle (3.4,3.4);
 \draw[manifoldred,very thick,-{Stealth[length=2.6mm]}]
   (1.0,0)--(2.2,0);
 \draw[manifoldred,very thick,-{Stealth[length=2.6mm]}]
   (1.0,3.4)--(2.2,3.4);
 \draw[manifoldgreen!80!black,very thick,-{Stealth[length=2.6mm]}]
   (0,1.0)--(0,2.2);
 \draw[manifoldgreen!80!black,very thick,-{Stealth[length=2.6mm]}]
   (3.4,1.0)--(3.4,2.2);
 \node[manifoldred,below=4pt] at (1.7,0) {$a$};
 \node[manifoldred,above=4pt] at (1.7,3.4) {$a$};
 \node[manifoldgreen!70!black,left=4pt] at (0,1.7) {$b$};
 \node[manifoldgreen!70!black,right=4pt] at (3.4,1.7) {$b$};
 \foreach \x in {0,3.4} \foreach \y in {0,3.4}
   \fill[manifoldgold!90!black] (\x,\y) circle(2.8pt);
 \node[below left=2pt] at (0,0) {$v$};
 \node[below right=2pt] at (3.4,0) {$v$};
 \node[above left=2pt] at (0,3.4) {$v$};
 \node[above right=2pt] at (3.4,3.4) {$v$};
 \node[manifoldblue!85!black] at (1.7,1.7) {$e^2$};
 \node at (1.7,-.83) {$\partial_2e^2=a+b-a-b=0$};
 % RP2: oba półokręgi obiegamy przeciwnie do wskazówek zegara.
 \node[font=\bfseries] at (8,4.1) {Płaszczyzna rzutowa $\mathbb{RP}^2$};
 \fill[manifoldblue!10] (8,1.7) circle(1.7);
 \draw[manifoldblue,very thick] (8,1.7) circle(1.7);
 \draw[manifoldred,very thick,-{Stealth[length=2.6mm]}]
   ({8+1.7*cos(45)},{1.7+1.7*sin(45)})
   arc[start angle=45,end angle=130,radius=1.7];
 \draw[manifoldred,very thick,-{Stealth[length=2.6mm]}]
   ({8+1.7*cos(225)},{1.7+1.7*sin(225)})
   arc[start angle=225,end angle=310,radius=1.7];
 \fill[manifoldgold!90!black] (6.3,1.7) circle(2.8pt)
   (9.7,1.7) circle(2.8pt);
 \node[left=5pt] at (6.3,1.7) {$v$};
 \node[right=5pt] at (9.7,1.7) {$v$};
 \node[manifoldred,above=5pt] at (8,3.4) {$a$};
 \node[manifoldred,below=5pt] at (8,0) {$a$};
 \node[manifoldblue!85!black] at (8,1.7) {$e^2$};
 \node at (8,-.83) {$\partial_2e^2=a+a=2a$};
\end{tikzpicture}
\caption{Po lewej sklejamy przeciwległe boki kwadratu zgodnie
ze strzałkami; wszystkie narożniki dają jeden punkt $v$.
Dodatni obieg brzegu czyta $aba^{-1}b^{-1}$, więc
oba kierunki występują z sumą współczynników zero.
Po prawej sklejamy antypody $z\sim-z$ na brzegu dysku.
Złote punkty utożsamiają się z $v$, a oba półokręgi
przechodzą na tę samą $1$-komórkę $a$ z tym samym zwrotem.
Pełny obieg brzegu daje dlatego współczynnik $2$.}
\label{fig:sklejenia-homologia-torus-rp2}
\end{figure}

\begin{figure}[!htbp]
\centering
\begin{tikzpicture}[>=Stealth,line cap=round,line join=round]
 \node[draw=manifoldblue!65!black,fill=manifoldblue!9,
  rounded corners=4pt,minimum width=2.15cm,
  minimum height=.8cm] (c2) at (1.2,0) {$C_2=\Z$};
 \node[draw=manifoldgreen!65!black,fill=manifoldgreen!9,
  rounded corners=4pt,minimum width=2.35cm,
  minimum height=.8cm] (c1) at (4.5,0) {$C_1=\Z^2$};
 \node[draw=manifoldgold!75!black,fill=manifoldgold!10,
  rounded corners=4pt,minimum width=2.15cm,
  minimum height=.8cm] (c0) at (7.8,0) {$C_0=\Z$};
 \draw[->,thick,manifoldred] (c2)--(c1)
  node[midway,above=5pt] {$\partial_2=0$};
 \draw[->,thick,manifoldred] (c1)--(c0)
  node[midway,above=5pt] {$\partial_1=0$};
 \node[align=center,font=\small,manifoldblue!75!black]
  at (1.2,-.9) {jedna\\$2$-komórka};
 \node[align=center,font=\small,manifoldgreen!75!black]
  at (4.5,-.9) {dwie\\pętle};
 \node[font=\small,manifoldgold!75!black]
  at (7.8,-.9) {punkt};
\end{tikzpicture}
\caption{Kompleks komórkowy torusa. Ma trzy
niezerowe grupy łańcuchów, a oba odwzorowania
brzegu znikają: pętle kończą się w tym
samym punkcie, a słowo dołączające
komórkę $2$-wymiarową jest komutatorem.}
\label{fig:kompleks-torusa}
\end{figure}

\begin{przyklad}[Sfera z dwóch komórek]
\label{prz:homologia-cw-sfera}
W modelu z
Przykładu~\ref{prz:cw-s2}
jest jedna komórka $e^0=v$
i jedna $e^2$, bez $1$-komórek.
Kompleks komórkowy to
\[
 0\longrightarrow\Z e^2
 \xrightarrow{\;0\;}0
 \xrightarrow{\;0\;}\Z v
 \longrightarrow0.
\]
Stąd $H_2(S^2)=\Z$, $H_1(S^2)=0$,
$H_0(S^2)=\Z$. Ogólniej
$S^n=D^n/\partial D^n$ ma
jedną komórkę w stopniu $0$
i jedną w stopniu $n$
($n\geq2$). Dla $S^1$
jedna $1$-komórka jest pętlą,
więc $\partial_1=0$.
W ten sposób odzyskujemy
Twierdzenie~\ref{thm:homologia-sfer}.
\end{przyklad}

\begin{przyklad}[Torus]
\label{prz:homologia-cw-torus}
Model CW z
Przykładu~\ref{prz:pi1-torus-rp2}
ma komórki $v$,
$a,b$ oraz $e^2$.
Krawędzie są pętlami, więc
$\partial_1a=v-v=0$ i tak samo $\partial_1b=0$.
Brzeg $e^2$ jest dołączony
według słowa $aba^{-1}b^{-1}$.
Aby obliczyć współczynnik przy $a$, zwijamy pętlę $b$ do punktu.
Pozostaje jeden dodatni i jeden ujemny obieg $a$, czyli
$\deg(q_a\circ\varphi)=1-1=0$. Po zwinięciu $a$ otrzymujemy
analogicznie $\deg(q_b\circ\varphi)=1-1=0$.
Ze Wzoru~\eqref{eq:formula-komorkowa} wynika więc $\partial_2e^2=0$.

Wyznaczmy teraz każdą grupę jako iloraz jądra przez obraz.
Nie ma $3$-komórek, więc $\operatorname{im}\partial_3=0$.
Ponieważ $\partial_2=0$, mamy
$\ker\partial_2=\Z e^2$ i $\operatorname{im}\partial_2=0$.
Podobnie $\ker\partial_1=\Z a\oplus\Z b$,
$\operatorname{im}\partial_1=0$ oraz $\ker\partial_0=\Z v$.
Stąd
\[
 \begin{aligned}
 H_2(T^2)&=\frac{\Z e^2}{0}\cong\Z,\\
 H_1(T^2)&=\frac{\Z a\oplus\Z b}{0}\cong\Z^2,\\
 H_0(T^2)&=\frac{\Z v}{0}\cong\Z.
 \end{aligned}
\]
W stopniach $k>2$ grupy łańcuchów są zerowe, więc $H_k(T^2)=0$.
Rysunek~\ref{fig:kompleks-torusa}
pokazuje odpowiadający kompleks.
Klasy $[a]$ i $[b]$ dają dwa niezależne kierunki obiegu torusa.
Klasa $[e^2]$ reprezentuje całą zorientowaną powierzchnię:
po sklejeniu przeciwległe boki znoszą się w jej brzegu.
To, że brzeg komórkowy jest zerowy, nie znaczy,
że odwzorowanie dołączające jest stałe. Komutator jest nietrywialny
w grupie podstawowej bukietu $S^1\vee S^1$, lecz
homologia zlicza sumy wykładników i daje tutaj zero.
\end{przyklad}

\begin{przyklad}[Płaszczyzna rzutowa i torsja]
\label{prz:homologia-rp2}
Przestrzeń $\mathbb{RP}^2$
ma po jednej komórce
$e^0=v$, $e^1=a$, $e^2$.
Jedyna $1$-komórka jest pętlą,
więc $\partial_1=0$.
Odwzorowanie dołączające $e^2$
okrąża $\mathbb{RP}^1\cong S^1$
dwa razy, co było widoczne
w Przykładzie~\ref{prz:pi1-torus-rp2}.
Ma stopień $2$ po ustaleniu
orientacji, zatem $\partial_2e^2=2a$ i kompleks ma postać
\[
 0\longrightarrow\Z e^2
 \xrightarrow{\;\times2\;}\Z a
 \xrightarrow{\;0\;}\Z v
 \longrightarrow0.
\]
Rachunek jąder i obrazów rozdziela dwa skutki tego samego
niezerowego operatora. Jeśli $\partial_2(ke^2)=2ka=0$,
to $2k=0$ w $\Z$, a więc $k=0$; stąd $\ker\partial_2=0$.
Nie ma zatem niezerowego $2$-cyklu komórkowego i $H_2=0$.
Natomiast $\ker\partial_1=\Z a$, a
$\operatorname{im}\partial_2=2\Z a$.
Łańcuch $ka$ jest brzegiem dokładnie wtedy, gdy $k$ jest parzyste.
Mamy więc
\[
 H_1(\mathbb{RP}^2)
 =\frac{\Z a}{2\Z a}\cong\Z/2\Z.
\]
Ponadto $H_0=\Z v/0\cong\Z$, a $H_k=0$ dla $k>2$.

Geometrycznie pojedynczy obieg $a$ nie jest brzegiem,
ale jego podwojenie jest brzegiem dołączonego dysku.
Klasa $[a]\ne0$ spełnia zatem $2[a]=0$.
To konkretny przykład \emph{torsji} zdefiniowanej w
Sekcji~\ref{sec:wspolczynniki-euler}.
Jednocześnie dysk sam nie daje całkowitoliczbowego cyklu powierzchni:
jego dwa półokręgi dodają się w $2a$, zamiast się znosić.
Rysunek~\ref{fig:sklejenia-homologia-torus-rp2} pokazuje,
dlaczego rachunek różni się od rachunku dla torusa.
\end{przyklad}

\begin{przyklad}[Dysk dołączony przez odwzorowanie stopnia $m$]
\label{prz:homologia-dolaczenie-stopien-m}
Niech $m\geq0$ i niech $f_m:S^1\to S^1$ będzie dane przez
$f_m(z)=z^m$ (dla $m=0$ jest to odwzorowanie stałe).
Do okręgu dołączmy dysk tym odwzorowaniem:
$M_m=S^1\cup_{f_m}D^2$.
Wybieramy jeden wierzchołek $v$ i jedną zorientowaną
$1$-komórkę $a$ w okręgu; wnętrze dysku to komórka $e^2$.
Stopień $f_m$ jest równy $m$, więc kompleks komórkowy to
\[
 0\longrightarrow\Z e^2
 \xrightarrow{\;\times m\;}\Z a
 \xrightarrow{\;0\;}\Z v\longrightarrow0.
\]
Dla $m\geq1$ mnożenie przez $m$ jest iniektywne,
czyli $\ker\partial_2=0$, a jego obraz to $m\Z a$.
Zatem $H_2(M_m)=0$, $H_1(M_m)=\Z/m\Z$ i $H_0(M_m)=\Z$;
wyższe homologie są zerowe.
Przy $m\geq2$ klasa $[a]$ ma dokładnie rząd $m$:
$ma=\partial_2e^2$, a dla $0<k<m$ liczba $k$ nie jest
wielokrotnością $m$, więc $ka$ nie jest brzegiem.
Dołączenie dysku zamienia tutaj nieskończony rząd obiegu okręgu
w skończony rząd jego klasy homologii.

Przypadki brzegowe warto rozpatrzyć osobno.
Dla $m=1$ odwzorowanie dołączające jest identycznością;
$M_1\cong D^2$, a operator $\partial_2$ jest izomorfizmem.
Stąd $H_1(M_1)=H_2(M_1)=0$, zgodnie ze ściągalnością dysku.
Dla $m=0$ cały brzeg dysku sklejamy z jednym punktem okręgu,
więc $M_0\cong S^1\vee S^2$.
Oba operatory brzegu są zerowe: teraz
$H_1(M_0)\cong\Z$ oraz $H_2(M_0)\cong\Z$.
Dołączana komórka nie usuwa żadnej klasy $H_1$, a tworzy nową klasę $H_2$.
Przypadek $m=2$ daje model $\mathbb{RP}^2$ z poprzedniego przykładu.
\end{przyklad}

\begin{uwaga}[Homologia ze współczynnikami]
\label{rem:homologia-wspolczynniki-przestrzeni}
Uzasadnijmy zmianę współczynników także w dodatniej charakterystyce.
Dla przemiennego pierścienia $R$ z jedynką definiujemy $C_n(X;R)$
jako wolny $R$-moduł na singularnych sympleksach, z tym samym wzorem brzegu.
Równoważnie jest to $C_n(X;\Z)\otimes_{\Z}R$.
Homotopie pryzmatyczne i podział barycentryczny mają całkowite współczynniki,
więc wszystkie ich tożsamości pozostają prawdziwe nad $R$.
Ciągi łańcuchów par oraz małych łańcuchów Mayera--Vietorisa są
rozszczepione w każdym stopniu na wybranych bazach, więc są dokładne także nad $R$.
Rachunek punktu, sfer i kolejnych szkieletów przebiega wtedy tak samo:
$C_n^{\mathrm{CW}}(X;R)$ ma po jednej kopii $R$ na komórkę.
Naturalność odwzorowań łączących względem $\Z\to R$ pokazuje,
że całkowity współczynnik brzegu $d$ staje się $d\cdot1_R$.
To uzasadnia obliczenia komórkowe nad $R$; nie zakładamy,
że utworzenie iloczynu tensorowego dowolnego ciągu dokładnego
z $R$ zawsze zachowuje dokładność.
\end{uwaga}

\begin{przyklad}[Te same przestrzenie nad różnymi ciałami]
\label{prz:zmiana-wspolczynnikow-rp2-torus}
Wstawmy teraz do powyższych kompleksów komórkowych
w miejsce $\Z$ dowolne
$\mathbb F\in\{\Q,\R,\C\}$.
Dla płaszczyzny rzutowej dostajemy
\[
 0\longrightarrow\mathbb F
 \xrightarrow{\;\times2\;}\mathbb F
 \xrightarrow{\;0\;}\mathbb F
 \longrightarrow0.
\]
Ponieważ $2$ ma w $\mathbb F$ odwrotność
$1/2$, środkowe odwzorowanie jest
izomorfizmem: jego jądro jest zerowe, a obraz równy całemu
$\mathbb F$. Dlatego $H_2=0$ oraz $H_1=\mathbb F/\mathbb F=0$.
Zatem
\[
 H_0(\mathbb{RP}^2;\mathbb F)=\mathbb F,
 \qquad H_k(\mathbb{RP}^2;\mathbb F)=0\quad(k>0).
\]
Klasa $\Z/2$ zniknęła. Dla torusa oba
operatory brzegu nadal są zerowe, więc
\[
 H_0(T^2;\mathbb F)=\mathbb F,\qquad
 H_1(T^2;\mathbb F)=\mathbb F^2,\qquad
 H_2(T^2;\mathbb F)=\mathbb F.
\]
Wymiary tych przestrzeni, odpowiednio
$1,2,1$, to liczby Betti'ego torusa.
Tak działa ogólny
Wzór~\eqref{eq:homologia-char-zero}.
Dla porównania nad $\Z/2$ strzałka
$\times2$ w kompleksie $\mathbb{RP}^2$
staje się zerowa, bo $2\cdot1=0$ w tym ciele.
Jej jądro jest teraz całym $\Z/2$, a obraz zerowy.
Wtedy oprócz $H_0$
otrzymujemy również
$H_1(\mathbb{RP}^2;\Z/2)\cong\Z/2$
oraz $H_2(\mathbb{RP}^2;\Z/2)\cong\Z/2$.
W rachunku modulo dwa podwójny brzeg $2a$ znika,
więc komórka $e^2$ staje się cyklem.
Nie uzyskaliśmy tego wyniku przez samo utworzenie iloczynu tensorowego
grup $H_k(\mathbb{RP}^2;\Z)$ z $\Z/2$: najpierw zmieniliśmy współczynniki
kompleksu, a dopiero potem policzyliśmy jego jądra i obrazy.
\end{przyklad}

\begin{przyklad}[Bukiet okręgów i charakterystyka Eulera]
\label{prz:homologia-bukiet}
Bukiet $r$ okręgów ma jeden
wierzchołek i $r$ pętli,
bez komórek wyższych:
$C_1^{\mathrm{CW}}=\Z^r$,
$C_0^{\mathrm{CW}}=\Z$ i
$\partial_1=0$.
Stąd $H_1=\Z^r$, $H_0=\Z$,
a wyższe homologie znikają.
Dla skończonego CW z
$c_n$ komórkami w wymiarze $n$
równanie Eulera z
Twierdzenia~\ref{thm:euler-homologia}
przyjmuje postać
\[
 \sum_n(-1)^n c_n
 =\sum_n(-1)^n
   \operatorname{rank}_{\Z} H_n(X;\Z).
\]
Wyjaśnijmy przejście do rang. Stosujemy twierdzenie Eulera
do kompleksu komórkowego nad $\Q$:
w stopniu $n$ ma on $c_n$ wektorów bazowych,
więc wymiar grupy łańcuchów wynosi $c_n$.
Z Twierdzenia~\ref{thm:homologia-char-zero}
$H_n(X;\Q)\cong H_n(X;\Z)\otimes_{\Z}\Q$,
a jego wymiar jest z Definicji~\ref{def:ranga-modulu}
równy $\operatorname{rank}_{\Z}H_n(X;\Z)$.
Podstawienie tych dwóch wymiarów daje wypisany wzór.
W szczególności dla torusa
$1-2+1=0$, a dla
$\mathbb{RP}^2$ mamy
$1-1+1=1$; torsja
$\Z/2\Z$ ma rangę zero.
\end{przyklad}


\section{Lokalna homologia i dalsza droga}
\label{sec:lokalna-homologia}

Grupy $H_k(X,X\setminus\{p\};\Z)$ nazywamy
\index{homologia lokalna@Homologia lokalna}
\emph{homologią lokalną przestrzeni $X$ w punkcie $p$}.
Łańcuchy omijające $p$ stają się w tej parze zerowe.
Pozostaje informacja o tym, jak przestrzeń otacza wyróżniony punkt.
W przestrzeni Hausdorffa wycinanie pozwoli obliczyć te grupy
w dowolnie małym otwartym otoczeniu $p$. Dla rozmaitości wynik zależy od
wymiaru oraz od tego, czy punkt leży we wnętrzu, czy na brzegu.

\begin{stwierdzenie}[Lokalna homologia rozmaitości]
\label{prop:lokalna-homologia}
Jeśli $M$ jest topologiczną
rozmaitością wymiaru $m$
bez brzegu i $p\in M$, to
\[
 H_k(M,M\setminus\{p\};\Z)
 \cong
 \begin{cases}\Z,&k=m,\\0,&k\ne m.\end{cases}
\]
Wybór generatora grupy w stopniu $m$
jest \emph{lokalną orientacją}
w punkcie $p$.
\end{stwierdzenie}
\begin{proof}
Wybierzmy otwarte otoczenie $U\subseteq M$ i mapę
$\psi:U\to B^m$ na otwartą kulę jednostkową,
z $\psi(p)=0$. Zbiór
$Z=M\setminus U$
jest domknięty i leży
we wnętrzu $M\setminus\{p\}$,
więc wycinanie
(Twierdzenie~\ref{thm:wycinanie})
daje
\[
 H_k(M,M\setminus\{p\})
 \cong H_k(U,U\setminus\{p\})
 \cong H_k(B^m,B^m\setminus\{0\}).
\]
Dla $m\geq1$ kula jest
ściągalna, a kula pozbawiona
środka silnie deformacyjnie retrahuje się na sferę
o promieniu $1/2$, homeomorficzną z $S^{m-1}$:
używamy $ru\mapsto((1-t)r+t/2)u$, co pozostaje w $0<r<1$.
Ponieważ $B^m\setminus\{0\}$ jest niepuste, zredukowany
długi ciąg pary daje
\[
 H_k(B^m,B^m\setminus\{0\})
 \cong\widetilde H_{k-1}(B^m\setminus\{0\})
 \cong\widetilde H_{k-1}(S^{m-1})
 \qquad(k\geq1).
\]
Twierdzenie~\ref{thm:homologia-sfer} wyznacza więc wszystkie
grupy dodatnich stopni. W stopniu zero inkluzja
$B^m\setminus\{0\}\hookrightarrow B^m$ jest surjektywna na $H_0$,
bo kula jest łukowo spójna i zbiór po lewej jest niepusty;
stąd $H_0(B^m,B^m\setminus\{0\})=0$.
Dla $m=0$ punkt jest otwartą składową
mapy lokalnej, a para lokalna to $(\{p\},\varnothing)$.
Jej grupa $H_0$ jest równa $\Z$, a wyższe grupy są zerowe.
\end{proof}

Istnienie lokalnego generatora
nie oznacza jeszcze, że można
wybrać generatory \emph{spójnie}
na całej rozmaitości. To będzie
treść pojęcia orientowalności. Najpierw porównajmy
otrzymany wynik z punktem brzegu rozmaitości
z Sekcji~\ref{sec:rozmaitosci-z-brzegiem-1}.

\begin{stwierdzenie}[Lokalna homologia w punkcie brzegu]
\label{prop:lokalna-homologia-brzeg-8}
Niech $M$ będzie topologiczną rozmaitością wymiaru $m\geq1$
z brzegiem i niech $p\in\partial M$.
Wtedy $H_k(M,M\setminus\{p\};\Z)=0$ dla każdego $k$.
\end{stwierdzenie}
\begin{proof}
Ograniczając mapę przy $p$, możemy znaleźć otwarte w $M$
otoczenie $U$ homeomorficzne z półkulą otwartą względem
półprzestrzeni, czyli ze zbiorem
\[
 B_+^m=\{x\in\R^m:\|x\|<1,\ x_m\geq0\},
\]
przy czym $p$ odpowiada $0$. Tak jak w poprzednim dowodzie,
wycinamy $M\setminus U$ i otrzymujemy
$H_k(M,M\setminus\{p\})\cong H_k(B_+^m,B_+^m\setminus\{0\})$.

Zbiór $B_+^m$ jest wypukły, więc jest ściągalny.
Dla $x=ru\ne0$, gdzie $0<r<1$, $\|u\|=1$ i $u_m\geq0$,
ten sam wzór radialny
$ru\mapsto((1-t)r+t/2)u$ daje silną retrakcję deformacyjną
$B_+^m\setminus\{0\}$ na
\[
 K=\{x\in\R^m:\|x\|=1/2,\ x_m\geq0\}.
\]
Jest to \emph{domknięta półsfera}, a nie cała sfera.
Rzutowanie na pierwsze $m-1$ współrzędnych daje homeomorfizm
$K\to\{y\in\R^{m-1}:\|y\|\leq1/2\}$,
którego odwrotnością jest
$y\mapsto(y,\sqrt{1/4-\|y\|^2})$.
Ostatnia współrzędna wynika z równania
$\|y\|^2+x_m^2=1/4$ i warunku $x_m\geq0$.
Zatem $K$ jest ściągalne; dla $m=1$ jest po prostu jednym punktem.
Wynika stąd, że $B_+^m\setminus\{0\}$ również jest niepuste
i ściągalne.

W zredukowanym długim ciągu tej pary obie grupy otaczające
$H_k(B_+^m,B_+^m\setminus\{0\})$ są zerowe dla $k\geq1$,
więc wszystkie te grupy względne znikają.
W stopniu zero inkluzja łączy dwie niepuste przestrzenie
łukowo spójne i indukuje izomorfizm $\Z\to\Z$.
Z dokładności otrzymujemy także zerowe $H_0$ pary.
\end{proof}

Różnica ma prostą przyczynę geometryczną: po usunięciu
środka kuli pozostaje sfera otaczająca brakujący punkt,
natomiast przy brzegu pozostaje tylko ściągalna półsfera.
Homologia lokalna rozróżnia więc punkt wewnętrzny i punkt
brzegu nawet wtedy, gdy cała rozmaitość jest ściągalna.
Na przykład odcinek $[0,1]$ ma homologię punktu,
lecz w jego punktach wewnętrznych lokalne $H_1$ jest równe $\Z$,
a na obu końcach wszystkie grupy lokalne są zerowe.

\begin{figure}[H]
\centering
\begin{tikzpicture}[line cap=round,line join=round,
 every node/.style={font=\small}]
 \begin{scope}
  \node[align=center] at (0,1.75) {punkt wewnętrzny\\odcinka};
  \draw[manifoldblue!70!black,line width=2pt] (-1.3,.5)--(1.3,.5);
  \fill[manifoldred] (0,.5) circle(3pt);
  \node[below=4pt] at (0,.5) {$p$};
  \node at (0,-.6) {$H_1\cong\Z$};
 \end{scope}
 \begin{scope}[xshift=4.4cm]
  \node[align=center] at (0,1.75) {punkt końcowy\\odcinka};
  \draw[manifoldblue!70!black,line width=2pt] (-1.05,.5)--(1.3,.5);
  \fill[manifoldred] (-1.05,.5) circle(3pt);
  \node[below=4pt] at (-1.05,.5) {$p$};
  \node at (0,-.6) {$H_1=0$};
 \end{scope}
 \begin{scope}[xshift=8.8cm]
  \node[align=center] at (0,1.75) {wierzchołek\\trójnogu};
  \draw[manifoldblue!70!black,line width=2pt]
   (-1.05,-.12)--(0,.5)--(1.05,-.12);
  \draw[manifoldblue!70!black,line width=2pt] (0,.5)--(0,1.22);
  \fill[manifoldred] (0,.5) circle(3pt);
  \node[right=4pt] at (0,.5) {$p$};
  \node at (0,-.6) {$H_1\cong\Z^2$};
 \end{scope}
\end{tikzpicture}
\caption{Porównanie grup $H_1(X,X\setminus\{p\};\Z)$.
Usunięcie wyróżnionego punktu pozostawia odpowiednio dwie,
jedną lub trzy składowe. Całe narysowane przestrzenie są
ściągalne, lecz mają różną homologię lokalną w punkcie $p$.}
\label{fig:lokalna-homologia-trzy-modele-8}
\end{figure}

\begin{przyklad}[Trójnóg nie jest rozmaitością w swoim wierzchołku]
\label{prz:lokalna-homologia-trojnog-8}
Niech $Y\subset\R^2$ będzie sumą trzech odcinków
$[0,u_1]$, $[0,u_2]$, $[0,u_3]$ leżących na trzech różnych
promieniach, a $p=0$ ich jedynym wspólnym punktem.
Przestrzeń $Y$ jest ściągalna przez homotopię
$(x,t)\mapsto(1-t)x$. Po usunięciu $p$ otrzymujemy trzy
składowe łukowe $(0,u_j]$, każdą ściągalną do swojego końca.
Zatem $H_0(Y\setminus\{p\})\cong\Z^3$, $H_0(Y)\cong\Z$,
a obie przestrzenie mają zerową homologię w stopniach dodatnich.

Zapiszmy potrzebny fragment długiego ciągu pary:
\[
 0\longrightarrow H_1(Y,Y\setminus\{p\})
 \xrightarrow{\ \delta\ }\Z^3
 \xrightarrow{\ s\ }\Z
 \longrightarrow H_0(Y,Y\setminus\{p\})
 \longrightarrow0.
\]
Wybierzmy generator $e_j$ grupy $H_0$ każdej składowej
$Y\setminus\{p\}$. W $Y$ wszystkie te punkty leżą w jednej
składowej, więc inkluzja posyła każdy $e_j$ na ten sam generator.
Stąd $s(a,b,c)=a+b+c$. Odwzorowanie $s$ jest surjektywne,
a jego jądro ma bazę $e_1-e_3$, $e_2-e_3$.
Istotnie, $a+b+c=0$ daje
$(a,b,c)=(a,b,-a-b)=a(e_1-e_3)+b(e_2-e_3)$;
współczynniki są jednoznaczne dzięki pierwszym dwóm współrzędnym.
Dokładność daje więc
\[
 H_1(Y,Y\setminus\{p\})\cong\ker s\cong\Z^2,
 \qquad H_0(Y,Y\setminus\{p\})=0.
\]
Pozostałe grupy lokalne również znikają: dla $k\geq2$
oba sąsiednie wyrazy $H_k(Y)$ i $H_{k-1}(Y\setminus\{p\})$
w długim ciągu są zerowe.

Gdyby $p$ miał otwarte otoczenie będące rozmaitością wymiaru $m$
z brzegiem, w punkcie wewnętrznym homologia lokalna byłaby równa
$\Z$ tylko w stopniu $m$, a w punkcie brzegu byłaby wszędzie zerowa.
Otrzymana grupa $H_1\cong\Z^2$ wyklucza każdy wymiar:
dla $m=1$ powinna być równa $\Z$, a dla $m\ne1$ zeru.
Problemem nie jest ściągalność
całego $Y$, lecz spotkanie trzech ramion w jednym punkcie.
\end{przyklad}



Na obecnym etapie mamy
już narzędzia do dalszej nauki:
kompleksy łańcuchowe i długie ciągi,
homotopijną niezmienniczość,
wycinanie, Mayera--Vietorisa
oraz model komórkowy, w którym
odwzorowania dołączające kontrolują
operator brzegu. Pozwalają one
później opisać, jak dołączenie
uchwytu zmienia homologię,
bez wprowadzania tutaj
teorii Morse'a i chirurgii.


\par\smallskip
{\footnotesize\noindent\emph{Dalsza lektura:}
Pełne rozwinięcie homologii symplicjalnej,
singularnej i komórkowej, wraz ze
szczegółami technicznymi dotyczącymi
podziału barycentrycznego, znajduje się
w Allen Hatcher,
\emph{Algebraic Topology}, rozdział 2,
\url{https://pi.math.cornell.edu/~hatcher/AT/AT.pdf}.\par}
\clearpage
